\appendices
\section{Complete Proofs for Sections III--VI}
\label{app:complete_proofs}

This appendix provides the complete proofs deferred from
Section~VI.  All notation and assumptions are inherited from
Sections~III--VI.  In particular,
$\boldsymbol\varpi^{\mathrm{tar}}$ is the public atomic deployment-target law,
and $\pi^{\mathrm{tar}}_{k,i}$ is the compatible client--atomic target law.
$P_r^{\mathrm{opp}}(k,i)$ is the opportunity-event law, while
$p^{\mathrm{obs}}_{k,r,i}$ and $q^{\mathrm{use}}_{k,r}$ are sequential
pre-outcome selection probabilities.  Finally, $\mathcal{C}_r$ is the feasible
aggregation set of (P2).

\subsection{Deferred Constructions from the Main Text}
\label{app:deferred_constructions}

The compatible client--stratum target used for opportunity correction is
\begin{equation}
    \Lambda_s^{\mathrm{tar}}
    =
    \sum_{i:s(i)=s}\varpi_i^{\mathrm{tar}}.
    \label{eq:target_stratum_mass}
\end{equation}
For $\Lambda_s^{\mathrm{tar}}>0$,
\begin{equation}
    \nu^{\mathrm{tar}}_{i\mid s}
    =
    \frac{\varpi_i^{\mathrm{tar}}}
         {\Lambda_s^{\mathrm{tar}}},
    \qquad s(i)=s,
    \label{eq:target_conditional_atomic}
\end{equation}
and $\lambda^{\mathrm{tar}}_{k\mid s}\ge0$ is a frozen public client share
with $\sum_k\lambda^{\mathrm{tar}}_{k\mid s}=1$.  Then
\begin{equation}
 \pi^{\mathrm{tar}}_{k,s}
 =
 \begin{cases}
 \Lambda_s^{\mathrm{tar}}\lambda^{\mathrm{tar}}_{k\mid s},
 & \Lambda_s^{\mathrm{tar}}>0,\\
 0, & \Lambda_s^{\mathrm{tar}}=0,
 \end{cases}
 \label{eq:target_ks_mass}
\end{equation}
and
\begin{equation}
 \pi^{\mathrm{tar}}_{k,i}
 =
 \begin{cases}
 \pi^{\mathrm{tar}}_{k,s(i)}
 \nu^{\mathrm{tar}}_{i\mid s(i)},
 & \Lambda_{s(i)}^{\mathrm{tar}}>0,\\
 0, & \Lambda_{s(i)}^{\mathrm{tar}}=0.
 \end{cases}
 \label{eq:target_ki_mass}
\end{equation}

\subsubsection{Target-Positive Pairs, Coverage, and the Heuristic Lookup}
\label{app:itar_lookup}
The set $I_{\mathrm{tar}}=\{(k,s):\pi^{\mathrm{tar}}_{k,s}>0\}$ and the
coverage $C_t$ in (6) use $\pi^{\mathrm{tar}}_{k,s}$ from
\eqref{eq:target_ks_mass}.  For unit $i$, $k(i)$ is the unique FL client
attached to $i$ on the frozen occupancy map: the spatial identity of $i$
under the same client map that builds \eqref{eq:target_ks_mass}.
SensorScope and U-Air attach stations; T-Drive and City Scanner attach
occupied cells.  The lookup is $\pi^{\mathrm{tar}}(k(i),s(i))$ on
that map; a missing key is $0$.  When several clients realize
opportunity on the same unit, the lookup still uses only the attached
client.  The descriptive coverage statistic of
Table~\ref{tab:scd_phi_freeze} is the train-unit
fraction of those zeros, not $\#\{i:\varpi_i^{\mathrm{tar}}=0\}$, and
not a positivity check on target-positive pairs.  It does not select the
Design factor of (32): that factor is configured, and
both of its settings are reported.
The pair opportunity mass $\pi^{\mathrm{opp}}_{k,s,t}$ is
$F_{t-}$-measurable on the same client--stratum key.
Because $\sum_{k,s}\pi^{\mathrm{tar}}_{k,s}=1$, (6) equals
the ratio of certified mass to total positive target mass.
Per-block lookup constructions and the published coverage statistic are in
Table~\ref{tab:scd_phi_freeze}.

A2's usable-update mean-ignorability display is
\begin{equation}
 \mathbb{E}\!\left[
 \left.
 \frac{U_{k,r}}{q^{\mathrm{use}}_{k,r}}
 \right|
 \mathcal{H}^{\mathrm{use}}_{k,r},
 O_e=1,k_e=k,i_e=i
 \right]
 =1.
 \label{eq:A2_use_mean_ignorability}
\end{equation}
The coordinate cap and $\ell_2$ support bound in (25) are
\begin{equation}
\begin{aligned}
 \overline\alpha_r
 &=
 \min\!\left\{
 1,
 \max\!\left\{
 \alpha_{\max},
 \frac{1}{K_r}
 \right\}
 \right\},\\
 \overline E_r
 &=
 \min\{E_{\min},K_r\}.
\end{aligned}
 \label{eq:cap_ess}
\end{equation}
The Lipschitz constant in (38) is
\begin{multline}
C_1
=
\frac{2K_r n_r}{B_{\min}}
\Bigl(
d_{\max}\max\{L_\pi,L_p\}+a_{\max}L_q
\Bigr),\\
L_\pi=\frac{\pi^{\mathrm{tar}}_{\max}}{\pi_{\min}^2 p_{\min}},\quad
L_p=\frac{\pi^{\mathrm{tar}}_{\max}}{\pi_{\min}p_{\min}^2},\\
L_q=\min\Bigl\{\frac{1}{q_{\min}^2},\,d_{\max}^2\Bigr\},
\label{eq:g4_c1}
\end{multline}
with $L_\pi=0$ and $L_p=1/p_{\min}^2$ under Design off.  These are the
constants of the implemented map: $L_\pi$ and $L_p$ bound its sensitivity to
the opportunity and observation margins, and $L_q$ its sensitivity to the
usable margin.  They are worst-case envelopes, not numerical certificates.
The ideal untruncated factorization of the sequential weight is
\begin{equation}
 \beta^\star_{k,r}\overline a^\star_{k,r,i}
 \propto
 R_{k,r,i}O_{k,r,i}U_{k,r}
 \frac{\zeta^{o}_{k,r,i}}
 {p^{\mathrm{obs}}_{k,r,i}q^{\mathrm{use}}_{k,r}},
 \label{eq:factorized_global_weight}
\end{equation}
where
$m^\star_{k,r}=\sum_i R_{k,r,i}O_{k,r,i}a^\star_{k,r,i}$,
$\overline a^\star_{k,r,i}=R_{k,r,i}O_{k,r,i}a^\star_{k,r,i}/m^\star_{k,r}$,
$d^\star_{k,r}=1/q^{\mathrm{use}}_{k,r}$, and
$\beta^\star_{k,r}=U_{k,r}m^\star_{k,r}d^\star_{k,r}/
\sum_j U_{j,r}m^\star_{j,r}d^\star_{j,r}$.

\subsection{Proof of Proposition~1}
\label{app:proof_prop1}

\begin{proposition*}[Target-vs-arrival risk distinction]
Let $\boldsymbol\varpi^{\mathrm{tar}}$ and
$\widetilde{\boldsymbol\rho}^{\mathrm{arr}}$ be two normalized laws on the same
finite supported atomic domain.  If
$\boldsymbol\varpi^{\mathrm{tar}}\neq
\widetilde{\boldsymbol\rho}^{\mathrm{arr}}$, then there exists a bounded
nonnegative error profile $e_i$ such that
\begin{equation}
    \sum_i \varpi_i^{\mathrm{tar}}e_i
    \neq
    \sum_i \widetilde\rho_i^{\mathrm{arr}}e_i.
    \label{eq:app_prop1_claim}
\end{equation}
Conversely, equality of the two weighted expectations for every bounded error
profile implies equality of the two underlying laws.
\end{proposition*}

\noindent\emph{Proof.}
Define $d_i\triangleq\varpi_i^{\mathrm{tar}}-\widetilde\rho_i^{\mathrm{arr}}$.
Because both vectors are probability laws,
\begin{equation}
    \sum_i d_i=0.
    \label{eq:app_prop1_zero_sum}
\end{equation}
The assumption
$\boldsymbol\varpi^{\mathrm{tar}}\neq
\widetilde{\boldsymbol\rho}^{\mathrm{arr}}$
implies that $d_i\neq0$ for at least one atomic unit.  Hence the set
$B_+\triangleq\{i:d_i>0\}$
is nonempty.  Indeed, if $d_i\le0$ for every $i$, then
\eqref{eq:app_prop1_zero_sum} would force $d_i=0$ for every $i$, contradicting
the assumed inequality of the two laws.

Choose the bounded nonnegative profile
$e_i=\mathbf{1}\{i\in B_+\}$, where $\mathbf{1}\{\cdot\}$ is the indicator of
the event in braces.
Then
\begin{align}
 &
 \sum_i
 \varpi_i^{\mathrm{tar}}e_i
 -
 \sum_i
 \widetilde\rho_i^{\mathrm{arr}}e_i
 \nonumber\\
 &\qquad=
 \sum_{i\in B_+}
 \left(
 \varpi_i^{\mathrm{tar}}
 -
 \widetilde\rho_i^{\mathrm{arr}}
 \right)
 =
 \sum_{i\in B_+}d_i
 >0.
 \label{eq:app_prop1_difference}
\end{align}
Therefore, the two weighted expectations differ for the bounded nonnegative
profile $e_i=\mathbf{1}\{i\in B_+\}$.  This proves the forward direction.

For the converse, suppose
\begin{equation}
    \sum_i \varpi_i^{\mathrm{tar}}e_i
    =
    \sum_i \widetilde\rho_i^{\mathrm{arr}}e_i
    \label{eq:app_prop1_all_equal}
\end{equation}
for every bounded error profile $e$.  Fix an arbitrary atomic unit $j$ and set
$e_i^{(j)}=\mathbf{1}\{i=j\}$.
Substituting $e_i^{(j)}$ into
\eqref{eq:app_prop1_all_equal} gives
$\varpi_j^{\mathrm{tar}}=\widetilde\rho_j^{\mathrm{arr}}$.
Since $j$ was arbitrary, the equality holds for every coordinate.  Hence
$\boldsymbol\varpi^{\mathrm{tar}}=\widetilde{\boldsymbol\rho}^{\mathrm{arr}}$.
This proves the converse. \hfill$\square$

\paragraph{Interpretation.}
The proposition proves that the two probability laws induce different linear
risk functionals whenever the laws differ.  It does not claim that a particular
fixed predictor must exhibit a nonzero numerical gap.  Nor does it claim that
the minimizers must differ for every restricted model class.


\subsection{Proof of Proposition~2}
\label{app:proof_prop2}

\begin{proposition*}[Normalized target reconstruction]
Under A1--A3, for any bounded test function $\psi(k,i)$, distinct from
the Design factor of (32), the H\'ajek ratio in
(35) converges in probability to
\begin{equation}
    \Phi^{\mathrm{tar}}
    =
    \sum_{k,i}
    \pi^{\mathrm{tar}}_{k,i}\psi(k,i).
    \label{eq:app_target_phi}
\end{equation}
\end{proposition*}

\noindent\emph{Proof.}
We first record two identities used by the proof.

\subsubsection{Target-Measure Compatibility}

For every stratum $s$ with $\Lambda_s^{\mathrm{tar}}>0$,
\eqref{eq:target_conditional_atomic} gives
\begin{equation}
    \sum_{i:s(i)=s}
    \nu^{\mathrm{tar}}_{i\mid s}
    =
    \frac{
    \sum_{i:s(i)=s}\varpi_i^{\mathrm{tar}}
    }{
    \Lambda_s^{\mathrm{tar}}
    }
    =1.
    \label{eq:app_nu_normalization}
\end{equation}
For an atomic unit $i$ whose stratum has positive target mass, using
$\pi^{\mathrm{tar}}_{k,s}
=\Lambda_s^{\mathrm{tar}}
\lambda^{\mathrm{tar}}_{k\mid s}$ and
$\sum_k\lambda^{\mathrm{tar}}_{k\mid s}=1$ gives
\begin{align}
 \sum_k\pi^{\mathrm{tar}}_{k,i}
 &=
 \sum_k
 \pi^{\mathrm{tar}}_{k,s(i)}
 \nu^{\mathrm{tar}}_{i\mid s(i)}
 \nonumber\\
 &=
 \Lambda_{s(i)}^{\mathrm{tar}}
 \nu^{\mathrm{tar}}_{i\mid s(i)}
 \sum_k
 \lambda^{\mathrm{tar}}_{k\mid s(i)}
 \nonumber\\
 &=
 \Lambda_{s(i)}^{\mathrm{tar}}
 \frac{\varpi_i^{\mathrm{tar}}}
 {\Lambda_{s(i)}^{\mathrm{tar}}}
 =
 \varpi_i^{\mathrm{tar}}.
 \label{eq:app_target_compatibility_positive}
\end{align}
If instead $\Lambda_{s(i)}^{\mathrm{tar}}=0$, then nonnegativity and
\eqref{eq:target_stratum_mass} imply $\varpi_i^{\mathrm{tar}}=0$.  The
zero-stratum convention in
\eqref{eq:target_ks_mass}--\eqref{eq:target_ki_mass}
gives $\pi^{\mathrm{tar}}_{k,i}=0$ for every $k$.  Hence, in both cases,
\begin{equation}
 \sum_k\pi^{\mathrm{tar}}_{k,i}
 =\varpi_i^{\mathrm{tar}}.
 \label{eq:app_target_compatibility}
\end{equation}
Summing over $i$ also yields
\begin{equation}
    \sum_{k,i}\pi^{\mathrm{tar}}_{k,i}
    =
    \sum_i\varpi_i^{\mathrm{tar}}
    =1.
    \label{eq:app_pi_normalized}
\end{equation}

\subsubsection{Sequential Inverse-Selection Identity}

Consider an opportunity event $e\in\mathcal{O}_r$ with mark
$(k_e,i_e)=(k,i)$.  For compactness, write
$p_e\triangleq p^{\mathrm{obs}}_{k,r,i}$ and
$q_{k,r}\triangleq q^{\mathrm{use}}_{k,r}$.
Under A1, both denominators are strictly positive on the target-supported event.

Let $\mathcal{H}^{\mathrm{obs}}_e$ denote the legal history immediately before
the observation outcome.  By the definition of $p_e$ and A2,
\begin{equation}
    \mathbb{E}
    \left[
    \left.
    \frac{O_e}{p_e}
    \right|
    k_e=k,\,
    i_e=i,\,
    \mathcal{H}^{\mathrm{obs}}_e
    \right]
    =1.
    \label{eq:app_obs_ipw}
\end{equation}
If $O_e=1$, client $k$ has at least one observed item and is therefore in the
registered-attempt set $\mathcal{E}_r$.  Conditional on the legal pre-$U$
history $\mathcal{H}^{\mathrm{use}}_{k,r}$, the definition of $q_{k,r}$ and A2
give
\begin{equation}
    \mathbb{E}
    \left[
    \left.
    \frac{U_{k,r}}{q_{k,r}}
    \right|
    k\in\mathcal{E}_r,\,
    \mathcal{H}^{\mathrm{use}}_{k,r}
    \right]
    =1.
    \label{eq:app_use_ipw}
\end{equation}
The usable-update indicator may be shared by several observed opportunity
events of the same client-window.  This induces within-window dependence, but
it does not invalidate the conditional identity
\eqref{eq:app_use_ipw}.  A3 handles such dependence at the
cluster/window level.

Applying the tower property in the temporal order
opportunity $\rightarrow$ observation $\rightarrow$ registration/use, and using
\eqref{eq:A2_use_mean_ignorability} on the $O_e=1$ branch, gives
\begin{align}
 &\mathbb{E}\!\left[
 \left.
 \frac{O_eU_{k,r}}{p_eq_{k,r}}
 \right|
 k_e=k,\,i_e=i,\,\mathcal{H}^{\mathrm{obs}}_e
 \right]
 \nonumber\\
 &\qquad=
 \mathbb{E}\!\left[
 \left.
 \frac{O_e}{p_e}
 \right|
 k_e=k,\,i_e=i,\,\mathcal{H}^{\mathrm{obs}}_e
 \right]
 =1.
 \label{eq:app_sequential_identity}
\end{align}
The inner conditional expectation is evaluated only on the $O_e=1$ branch,
where client $k$ is registered.  It equals one by the strengthened
mean-ignorability condition \eqref{eq:A2_use_mean_ignorability}.  The outer
expectation is then \eqref{eq:app_obs_ipw}.

With these two identities in hand, we prove the proposition.
For event $e\in\mathcal{O}_r$, recall
\begin{equation}
 \chi_e
 =
 O_eU_{k_e,r}
 \frac{
 \pi^{\mathrm{tar}}_{k_e,i_e}
 }{
 P_r^{\mathrm{opp}}(k_e,i_e)
 p^{\mathrm{obs}}_{k_e,r,i_e}
 q^{\mathrm{use}}_{k_e,r}
 }.
 \label{eq:app_chi}
\end{equation}
Condition first on the opportunity mark $(k_e,i_e)=(k,i)$.  Using
\eqref{eq:app_sequential_identity},
\begin{align}
 &\mathbb{E}
 \left[
 \left.
 \chi_e\psi(k_e,i_e)
 \right|
 k_e=k,\,
 i_e=i
 \right]
 \nonumber\\
 &=
 \frac{
 \pi^{\mathrm{tar}}_{k,i}
 }{
 P_r^{\mathrm{opp}}(k,i)
 }
 \psi(k,i)
 \,
 \mathbb{E}
 \left[
 \left.
 \frac{
 O_eU_{k,r}
 }{
 p^{\mathrm{obs}}_{k,r,i}
 q^{\mathrm{use}}_{k,r}
 }
 \right|
 k_e=k,\,
 i_e=i
 \right]
 \nonumber\\
 &=
 \frac{
 \pi^{\mathrm{tar}}_{k,i}
 }{
 P_r^{\mathrm{opp}}(k,i)
 }
 \psi(k,i).
 \label{eq:app_conditional_weighted_moment}
\end{align}
By (1), the mark itself has opportunity law
$P_r^{\mathrm{opp}}(k,i)$.  Averaging
\eqref{eq:app_conditional_weighted_moment} over this law therefore yields
\begin{align}
 \mathbb{E}
 \left[
 \chi_e\psi(k_e,i_e)
 \right]
 &=
 \sum_{k,i}
 P_r^{\mathrm{opp}}(k,i)
 \frac{
 \pi^{\mathrm{tar}}_{k,i}
 }{
 P_r^{\mathrm{opp}}(k,i)
 }
 \psi(k,i)
 \nonumber\\
 &=
 \sum_{k,i}
 \pi^{\mathrm{tar}}_{k,i}\psi(k,i)
 =
 \Phi^{\mathrm{tar}}.
 \label{eq:app_expected_numerator_event}
\end{align}
Likewise, taking $\psi\equiv1$ and using
\eqref{eq:app_pi_normalized},
\begin{equation}
    \mathbb{E}[\chi_e]
    =
    \sum_{k,i}\pi^{\mathrm{tar}}_{k,i}
    =1.
    \label{eq:app_expected_denominator_event}
\end{equation}

Let
\begin{equation}
    N_R^{\mathrm{opp}}
    \triangleq
    \sum_{r=0}^{R-1}|\mathcal{O}_r|
    \label{eq:app_total_opp_events}
\end{equation}
be the total number of opportunity events in the first $R$ windows.  Define
\begin{align}
 A_R(\psi)
 &\triangleq
 \sum_{r<R}
 \sum_{e\in\mathcal{O}_r}
 \chi_e\psi(k_e,i_e),
 \label{eq:app_AR}\\
 B_R
 &\triangleq
 \sum_{r<R}
 \sum_{e\in\mathcal{O}_r}
 \chi_e.
 \label{eq:app_BR}
\end{align}
A3 explicitly requires the asymptotic regime
$N_R^{\mathrm{opp}}\rightarrow\infty$.  It is precisely the regularity
condition needed to pass from the eventwise identities above to the possibly
dependent client-window process.  Along this growing sequence, it gives
\begin{align}
 \frac{A_R(\psi)}{N_R^{\mathrm{opp}}}
 &\xrightarrow{p}
 \Phi^{\mathrm{tar}},
 \label{eq:app_AR_lln}\\
 \frac{B_R}{N_R^{\mathrm{opp}}}
 &\xrightarrow{p}
 1.
 \label{eq:app_BR_lln}
\end{align}
The second limit is strictly positive.  Therefore, by the continuous mapping
theorem (equivalently, Slutsky's theorem),
\begin{align}
 \widehat\Phi_R
 &=
 \frac{A_R(\psi)}{B_R}
 =
 \frac{
 A_R(\psi)/N_R^{\mathrm{opp}}
 }{
 B_R/N_R^{\mathrm{opp}}
 }
 \xrightarrow{p}
 \Phi^{\mathrm{tar}}.
 \label{eq:app_hajek_limit}
\end{align}
This proves the main convergence statement.

It remains to verify the two stated marginals.

First, fix $(k_0,s_0)$ and choose
\begin{equation}
    \psi_{k_0,s_0}(k,i)
    =
    \mathbf{1}\{k=k_0,\ s(i)=s_0\}.
    \label{eq:app_phi_client_stratum}
\end{equation}
If $\Lambda_{s_0}^{\mathrm{tar}}>0$, then
\begin{align}
 \Phi^{\mathrm{tar}}
 &=
 \sum_{i:s(i)=s_0}
 \pi^{\mathrm{tar}}_{k_0,i}
 \nonumber\\
 &=
 \pi^{\mathrm{tar}}_{k_0,s_0}
 \sum_{i:s(i)=s_0}
 \nu^{\mathrm{tar}}_{i\mid s_0}
 =
 \pi^{\mathrm{tar}}_{k_0,s_0},
 \label{eq:app_recover_client_stratum_positive}
\end{align}
where the last equality follows from
\eqref{eq:app_nu_normalization}.  If
$\Lambda_{s_0}^{\mathrm{tar}}=0$, the zero-stratum convention gives both
$\pi^{\mathrm{tar}}_{k_0,s_0}=0$ and
$\pi^{\mathrm{tar}}_{k_0,i}=0$ for every $i$ in that stratum.  The same
conclusion then holds.  Thus, in all cases,
$\Phi^{\mathrm{tar}}=\pi^{\mathrm{tar}}_{k_0,s_0}$.

Second, fix a deployment group $g$ and choose
$\psi_g(k,i)=\mathbf{1}\{h(i)=g\}$.
Using target compatibility
\eqref{eq:app_target_compatibility},
\begin{align}
 \Phi^{\mathrm{tar}}
 &=
 \sum_{k,i}
 \pi^{\mathrm{tar}}_{k,i}
 \mathbf{1}\{h(i)=g\}
 \nonumber\\
 &=
 \sum_{i:h(i)=g}
 \sum_k
 \pi^{\mathrm{tar}}_{k,i}
 \nonumber\\
 &=
 \sum_{i:h(i)=g}
 \varpi_i^{\mathrm{tar}}
 =
 \mu_g.
 \label{eq:app_recover_group}
\end{align}
Hence, the ideal normalized sequential correction recovers both the
client--stratum target moments and, through the compatibility bridge, the
deployment-group mass.  It does not prove optimization of
$\mathcal{R}_{\mu}(\theta)$. \hfill$\square$

\paragraph{Scope of Proposition~2.}
The result is an identification/consistency statement for the ideal pooled
sequential ratios: compatible target moments or target-group mass under
A1--A3.  It does not imply that every finite-window normalized
reference $\boldsymbol\beta_r^\star$ exactly realizes the deployment target,
and it does not prove deployment-target risk optimization, a $Z\!\to\!Y$
outcome model, unbiased local gradients, or convergence under asynchrony
and staleness.
The current mobile support can make exact realization infeasible, even when
the pooled sequential weights have the correct target-law interpretation.  This
is the role of the feasibility-aware calibration in
Section~V.  The result also does not assert exact
unbiasedness of the finite-sample clipped plug-in implementation.  When the
stabilization bounds are active, the practical estimator deliberately trades
exact inverse-weight identification for controlled weight concentration.


\subsection{Proof of Theorem~1}
\label{app:proof_theorem1}

\noindent\textbf{Theorem~1.}
Under A4, every active window has a nonempty compact convex feasible set
$\mathcal{C}_r$, and the calibration program
$\min_{\boldsymbol\alpha\in\mathcal{C}_r}J_r(\boldsymbol\alpha)$
has a unique global minimizer $\boldsymbol\alpha_r^\star$.

\noindent\emph{Proof.}
Fix an active window $r$.  Then
\begin{equation}
    K_r=|\mathcal{A}_r|\ge1.
    \label{eq:app_K_positive}
\end{equation}
Consider the uniform vector
\begin{equation}
    \boldsymbol\alpha_r^{u}
    =
    \frac{1}{K_r}\mathbf{1}.
    \label{eq:app_uniform_alpha}
\end{equation}
It satisfies
\begin{equation}
    \mathbf{1}^{\top}\boldsymbol\alpha_r^{u}=1
    \quad\text{and}\quad
    \alpha_{k,r}^{u}=\frac{1}{K_r}\ge0.
    \label{eq:app_uniform_simplex}
\end{equation}
From \eqref{eq:cap_ess},
\begin{equation}
    \overline\alpha_r
    =
    \min\left\{
    1,
    \max\left\{
    \alpha_{\max},
    \frac{1}{K_r}
    \right\}
    \right\}.
    \label{eq:app_alpha_bar}
\end{equation}
Since $K_r\ge1$ and $0<\alpha_{\max}\le1$,
\begin{equation}
    \frac{1}{K_r}\le\overline\alpha_r\le1.
    \label{eq:app_uniform_box}
\end{equation}
Thus the box constraints are satisfied by
$\boldsymbol\alpha_r^{u}$.

Next,
\begin{equation}
    \left\|
    \boldsymbol\alpha_r^{u}
    \right\|_2^2
    =
    K_r\left(\frac{1}{K_r}\right)^2
    =
    \frac{1}{K_r}.
    \label{eq:app_uniform_l2}
\end{equation}
Because
\begin{equation}
    \overline E_r
    =
    \min\{E_{\min},K_r\}
    \le K_r,
    \label{eq:app_Ebar_leq_K}
\end{equation}
we have
\begin{equation}
    \frac{1}{K_r}
    \le
    \frac{1}{\overline E_r}.
    \label{eq:app_uniform_ess}
\end{equation}
Therefore
\begin{equation}
    \boldsymbol\alpha_r^{u}
    \in
    \mathcal{C}_r,
    \label{eq:app_C_nonempty}
\end{equation}
so $\mathcal{C}_r$ is nonempty.

The equality constraint
$\mathbf{1}^{\top}\boldsymbol\alpha=1$ defines an affine set.  The coordinate
constraints
$0\le\alpha_k\le\overline\alpha_r$ define a closed convex box, and
\begin{equation}
    \left\{
    \boldsymbol\alpha:
    \|\boldsymbol\alpha\|_2^2
    \le
    1/\overline E_r
    \right\}
    \label{eq:app_l2_ball}
\end{equation}
is a closed convex Euclidean ball.  Their intersection is therefore closed and
convex.  Moreover, the simplex and box constraints imply
$0\le\alpha_k\le1$ for every coordinate, so the set is bounded.  Hence
$\mathcal{C}_r$ is compact.

We now analyze the objective.  Expanding only the terms that are quadratic in
$\boldsymbol\alpha$, the Hessian of $J_r$ is
\begin{equation}
 \nabla^2J_r
 =
 \lambda_gM_r^{\top}M_r
 +
 \lambda_{\beta}I
 +
 \lambda_vV_r.
 \label{eq:app_P2_hessian}
\end{equation}
The debt term
$-\mathbf{Q}_r^{\top}M_r\boldsymbol\alpha$ and the staleness term
$\lambda_s\overline{\boldsymbol\tau}_r^{\top}\boldsymbol\alpha$ are linear.
They therefore contribute zero Hessian.  Under A4,
\begin{equation}
    M_r^{\top}M_r\succeq0,
    \qquad
    V_r\succeq0,
    \qquad
    \lambda_g,\lambda_v\ge0,
    \qquad
    \lambda_{\beta}>0.
    \label{eq:app_psd_terms}
\end{equation}
Consequently, for every nonzero vector $\mathbf{x}\in\mathbb{R}^{K_r}$,
\begin{align}
 \mathbf{x}^{\top}\nabla^2J_r\mathbf{x}
 &=
 \lambda_g\|M_r\mathbf{x}\|_2^2
 +
 \lambda_{\beta}\|\mathbf{x}\|_2^2
 +
 \lambda_v\mathbf{x}^{\top}V_r\mathbf{x}
 \nonumber\\
 &\ge
 \lambda_{\beta}\|\mathbf{x}\|_2^2
 >0.
 \label{eq:app_strong_convex}
\end{align}
Thus $J_r$ is $\lambda_{\beta}$-strongly convex on
$\mathbb{R}^{K_r}$ and hence strictly convex on $\mathcal{C}_r$.

Because $J_r$ is continuous and $\mathcal{C}_r$ is nonempty and compact, the
Weierstrass theorem guarantees the existence of at least one minimizer.
Strict convexity on the convex feasible set guarantees that two distinct
minimizers cannot exist.  Therefore, (P2) has exactly one global minimizer,
denoted by $\boldsymbol\alpha_r^\star$.  This completes the proof.
\hfill$\square$

\paragraph{Scope of Theorem~1.}
The theorem proves existence and uniqueness of the mathematical optimizer of
(P2).  It does not, by itself, certify that a finite-precision solver log
reconstructs that optimizer exactly.  Numerical feasibility and post-hoc
solver audit are separate empirical/provenance questions.


\subsection{Proof of Proposition~3}
\label{app:proof_prop3}

\noindent\textbf{Proposition~3.}
Assume A1, A4, A5, and A6b, let
$\mathbf{1}^{\top}b_r^{\mathcal T}\ge B_{\min}>0$ hold at the true margins,
and let the ideal and the implemented reference share $\mathcal{A}_r$,
$M_r$, and $\mathcal{C}_r$.  Then the implemented clipped plug-in
$\widehat{\boldsymbol\beta}_r$ is a Lipschitz perturbation of the
same-window ideal reference $\boldsymbol\beta_r^\star$ up to the stabilization
error $\varepsilon_{\mathrm{stab},r}$, and the P2 map does not amplify that
$\ell_2$ error on a fixed feasible set.

\noindent\emph{Proof.}
Fix an active window $r$ and the common set $\mathcal{A}_r$.  Write
$\widehat a$ and $\widehat d$ for the implemented plug-in ratios
(16) and (23) evaluated at the estimates, and
$a^{\mathcal T}$ and $d^{\mathcal T}$ for the same two maps evaluated at the
true margins,
\[
\begin{aligned}
a^{\mathcal T}_{k,r,i}
&=
\min\Bigl\{a_{\max},
\frac{\zeta_{k,r,i}}{\max\{p^{\mathrm{obs}}_{k,r,i},p_{\min}\}}\Bigr\},\\
d^{\mathcal T}_{k,r}
&=
\min\Bigl\{d_{\max},
\frac{1}{\max\{q^{\mathrm{use}}_{k,r},q_{\min}\}}\Bigr\},
\end{aligned}
\]
where $\zeta=\pi^{\mathrm{tar}}/\pi^{\mathrm{opp}}$ under Design on and
$\zeta=1$ under Design off, and write $m^{\mathcal T}$, $b^{\mathcal T}$ and
$\boldsymbol\beta_r^{\mathcal T}$ for the corresponding masses and normalized
reference.

\emph{Step 1 (the stabilization map is Lipschitz without A6a).}  The floors
and ceilings are the only sources of non-smoothness: $x\mapsto\max\{x,c\}$ and
$x\mapsto\min\{a_{\max},x\}$ are non-expansive and $x\mapsto1/x$ is
$1/c^2$-Lipschitz on $[c,\infty)$, so $x\mapsto1/\max\{x,c\}$ is
$1/c^2$-Lipschitz on $(0,\infty)$.  With $\pi^{\mathrm{tar}}\le1$,
\begin{equation}
\begin{aligned}
 \bigl\lvert a^{\mathcal T}-\widehat a\bigr\rvert
 &\le
 L_\pi\bigl\lvert\pi^{\mathrm{opp}}-\widehat\pi^{\mathrm{opp}}\bigr\rvert\\
 &\quad
 +L_p\bigl\lvert p^{\mathrm{obs}}-\widehat p^{\mathrm{obs}}\bigr\rvert,\\
 \bigl\lvert d^{\mathcal T}-\widehat d\bigr\rvert
 &\le
 L_q\bigl\lvert q^{\mathrm{use}}-\widehat q^{\mathrm{use}}\bigr\rvert,
\end{aligned}
 \label{eq:app_plugin_lipschitz}
\end{equation}
with the constants of \eqref{eq:g4_c1}.  Under Design off the opportunity
margin does not enter the map at all, so $\zeta\equiv1$ gives $L_\pi=0$ and
$L_p=1/p_{\min}^2$.  A6b bounds the right-hand sides by
$L_\pi\varepsilon_\pi+L_p\varepsilon_p$ and $L_q\varepsilon_q$.

\emph{Step 2 (mass error).}  Each sum in (17) has at most $n_r$
terms, so
$\lvert m_k-m^{\mathcal T}_{k}\rvert
\le n_r\max_i\lvert a_{k,r,i}-a^{\mathcal T}_{k,r,i}\rvert$, and with
$b_k=m_kd_{k,r}$,
\begin{equation}
\begin{aligned}
 \bigl\lvert b_k-b^{\mathcal T}_k\bigr\rvert
 \le{}&
 d_{\max}\,n_r\,\bigl\|a-a^{\mathcal T}\bigr\|_\infty\\
 &+a_{\max}\,n_r\,\bigl\|d-d^{\mathcal T}\bigr\|_\infty .
\end{aligned}
 \label{eq:app_plugin_mass}
\end{equation}
Only the implementation-side bounds $a\le a_{\max}$ and $d\le d_{\max}$ enter
here; the ideal ratios need not lie below the ceilings, so this step does not
need A6a either.

\emph{Step 3 (H\'ajek normalization).}  For nonnegative $x,y$ with
$\mathbf{1}^{\top}x>0$,
\[
\Bigl\|\frac{x}{\mathbf{1}^{\top}x}
-\frac{y}{\mathbf{1}^{\top}y}\Bigr\|_1
\le\frac{2\|x-y\|_1}{\mathbf{1}^{\top}x},
\]
by inserting $\pm y/\mathbf{1}^{\top}x$.  Applied with $x=b^{\mathcal T}$,
whose mass is bounded below by $B_{\min}$, and combined with Step 2,
\[
\bigl\|\widehat{\boldsymbol\beta}_r
-\boldsymbol\beta_r^{\mathcal T}\bigr\|_1
\le
\frac{2K_rn_r}{B_{\min}}
\Bigl(
d_{\max}\bigl\|a-a^{\mathcal T}\bigr\|_\infty
+a_{\max}\bigl\|d-d^{\mathcal T}\bigr\|_\infty
\Bigr).
\]

\emph{Step 4 (triangle inequality).}  By (37),
$\|\widehat{\boldsymbol\beta}_r-\boldsymbol\beta_r^\star\|_1
\le\|\widehat{\boldsymbol\beta}_r-\boldsymbol\beta_r^{\mathcal T}\|_1
+\varepsilon_{\mathrm{stab},r}$, and Steps 1--3 bound the first term by
$C_1(\varepsilon_\pi+\varepsilon_p+\varepsilon_q)$ with the $C_1$ of
\eqref{eq:g4_c1}: that is (38).

\emph{Support lemma.}  Under Design on,
$a>0\iff\pi^{\mathrm{tar}}>0\iff a^\star>0$ on A5's active support; under
Design off every ratio is strictly positive there.  Hence $m^{\mathcal T}$,
$b^{\mathcal T}$, $m^\star$ and $b^\star$ share a support, which is the
shared-support clause of A5.

The P2 bound (39) is the already-proved reference
non-expansion \eqref{eq:app_reference_nonexpansive} on the common
convex set $\mathcal{C}_r$.  The composition bound
(40) is the operator $2$-norm.
\hfill$\square$

The constant $C_1$ is a worst-case envelope, not a numerical certificate: at
the frozen constants its Design-on value is vacuous by many orders of
magnitude, so the informative term is $\varepsilon_{\mathrm{stab},r}$.


\subsection{Proof of the Target-Debt Accounting Bound}
\label{app:proof_debt_bound}

The bound states that the prefix-averaged deployment composition
$\overline{\boldsymbol\omega}_R$ cannot drift farther from the target
$\boldsymbol\mu$ than twice the accumulated target debt per unit of active
step mass.  Intuitively, the debt vector records, for each deployment group,
how much service the group has cumulatively missed across windows.  The
relation therefore turns that accounting into a deterministic control on
structural mismatch: the group-average distance to the target is bounded by
how much the process has fallen behind, normalized by how much it has run.
The proof proceeds in three steps.  It first shows that the coordinate-wise
debt dominates the cumulative signed discrepancy.  It then combines this
domination with the zero-sum property of the discrepancy to bound its
$\ell_1$ norm by twice the debt.  Finally, it divides by the active step
mass $S_R$ to obtain the stated bound.

We now prove the deterministic relation in
(27).  Recall
\begin{equation}
    S_R
    =
    \sum_{r=0}^{R-1}\eta_rI_r,
    \qquad
    \overline{\boldsymbol\omega}_R
    =
    \frac{
    \sum_{r=0}^{R-1}
    \eta_rI_r\boldsymbol\omega_r
    }{S_R},
    \label{eq:app_debt_prefix_defs}
\end{equation}
for $S_R>0$.  In every active window,
$\boldsymbol\alpha_r$ is a probability vector and each column
$\mathbf{c}_{k,r}^{g}$ of $M_r$ is a deployment-group probability vector.
Therefore
\begin{equation}
    \boldsymbol\omega_r
    =
    M_r\boldsymbol\alpha_r
    \in
    \Delta^{|\mathcal G|-1},
    \label{eq:app_omega_simplex}
\end{equation}
and hence
\begin{equation}
    \mathbf{1}^{\top}\boldsymbol\omega_r
    =
    \mathbf{1}^{\top}\boldsymbol\mu
    =1.
    \label{eq:app_equal_mass}
\end{equation}

Define the cumulative signed target discrepancy
\begin{equation}
    \mathbf{D}_R
    \triangleq
    \sum_{r=0}^{R-1}
    \eta_rI_r
    \left(
    \boldsymbol\mu-\boldsymbol\omega_r
    \right).
    \label{eq:app_DR}
\end{equation}
By \eqref{eq:app_debt_prefix_defs},
\begin{equation}
    \mathbf{D}_R
    =
    S_R
    \left(
    \boldsymbol\mu-\overline{\boldsymbol\omega}_R
    \right).
    \label{eq:app_DR_baromega}
\end{equation}

Fix a target-group coordinate $g$.  In an active window, the debt recursion is
\begin{equation}
    Q_{r+1,g}
    =
    \left[
    Q_{r,g}
    +
    \eta_r
    \left(
    \mu_g-\omega_{r,g}
    \right)
    \right]_+.
    \label{eq:app_debt_coordinate}
\end{equation}
Since $[x]_+\ge x$ for every scalar $x$,
\begin{equation}
    Q_{r+1,g}
    \ge
    Q_{r,g}
    +
    \eta_r
    \left(
    \mu_g-\omega_{r,g}
    \right)
    \label{eq:app_debt_one_step}
\end{equation}
holds for active windows.  In inactive windows,
$Q_{r+1,g}=Q_{r,g}$ and $I_r=0$.  The same inequality can therefore be written
uniformly as
\begin{equation}
    Q_{r+1,g}
    \ge
    Q_{r,g}
    +
    \eta_rI_r
    \left(
    \mu_g-\omega_{r,g}
    \right).
    \label{eq:app_debt_one_step_all}
\end{equation}
Telescoping from $r=0$ to $R-1$ and using $Q_{0,g}=0$ gives
\begin{equation}
    Q_{R,g}
    \ge
    D_{R,g}.
    \label{eq:app_Q_dominates_D}
\end{equation}
Because $Q_{R,g}\ge0$, it follows that
\begin{equation}
    (D_{R,g})_+
    \le
    Q_{R,g}.
    \label{eq:app_positive_part_domination}
\end{equation}
Summing over $g$,
\begin{equation}
    \|(\mathbf{D}_R)_+\|_1
    \le
    \|\mathbf{Q}_R\|_1.
    \label{eq:app_positive_part_l1}
\end{equation}

Next, by \eqref{eq:app_equal_mass},
\begin{align}
    \mathbf{1}^{\top}\mathbf{D}_R
    &=
    \sum_{r=0}^{R-1}
    \eta_rI_r
    \mathbf{1}^{\top}
    \left(
    \boldsymbol\mu-\boldsymbol\omega_r
    \right)
    \nonumber\\
    &=
    \sum_{r=0}^{R-1}
    \eta_rI_r(1-1)
    =0.
    \label{eq:app_DR_zero_sum}
\end{align}
For any finite-dimensional vector with zero coordinate sum, the positive and
negative parts have equal $\ell_1$ mass:
\begin{equation}
    \|(\mathbf{D}_R)_+\|_1
    =
    \|(\mathbf{D}_R)_-\|_1,
    \label{eq:app_pos_neg_equal}
\end{equation}
where $(\mathbf{x})_-=[-\mathbf{x}]_+$.  Therefore
\begin{equation}
\begin{aligned}
    \|\mathbf{D}_R\|_1
    &=
    \|(\mathbf{D}_R)_+\|_1
    +
    \|(\mathbf{D}_R)_-\|_1\\
    &=2\|(\mathbf{D}_R)_+\|_1
    \le
    2\|\mathbf{Q}_R\|_1.
\end{aligned}
    \label{eq:app_DR_l1_bound}
\end{equation}
Finally, substituting
\eqref{eq:app_DR_baromega} into
\eqref{eq:app_DR_l1_bound} and dividing by $S_R>0$ yields
\begin{equation}
    \left\|
    \overline{\boldsymbol\omega}_R
    -
    \boldsymbol\mu
    \right\|_1
    \le
    \frac{
    2\|\mathbf{Q}_R\|_1
    }{
    S_R
    }.
    \label{eq:app_debt_bound_final}
\end{equation}
This is exactly the accounting relation stated in
(27). \hfill$\square$

\paragraph{Consequence.}
If the cumulative debt is sublinear in active step mass, that is,
$\|\mathbf{Q}_R\|_1=o(S_R)$, then
\begin{equation}
    \Delta_{\mathrm{group}}(R)
    =
    \left\|
    \overline{\boldsymbol\omega}_R
    -
    \boldsymbol\mu
    \right\|_1
    \longrightarrow0.
    \label{eq:app_debt_sublinear_consequence}
\end{equation}
This is a conditional implication of the deterministic bound.  It is not a
separate claim that RAVEN guarantees sublinear debt under arbitrary mobile
opportunity processes.


\subsection{Supplementary Consequences of the Calibration Geometry}
\label{app:supp_calibration_geometry}

The following two properties are deterministic consequences of the feasible
set and objective defined in Section~V.  They make the
finite-window geometry explicit; they are not additional main-paper theorem
claims.

\paragraph{Feasible target image.}
For an active window, define the set of deployment-group compositions that can
be realized by the currently available clients as
\begin{equation}
 \Omega_r
 \triangleq
 \left\{
 M_r\boldsymbol\alpha:
 \boldsymbol\alpha\in\mathcal C_r
 \right\}.
 \label{eq:app_feasible_target_image}
\end{equation}
Under A4, $\mathcal C_r$ is nonempty, compact, and convex.  Since
$\boldsymbol\alpha\mapsto M_r\boldsymbol\alpha$ is linear, $\Omega_r$ is also
nonempty, compact, and convex.  Hence, the irreducible instantaneous
feasibility gap
\begin{equation}
 d_r^{\mathrm{feas}}
 \triangleq
 \min_{\boldsymbol\omega\in\Omega_r}
 \left\|
 \boldsymbol\omega-\boldsymbol\mu
 \right\|_2
 \label{eq:app_feasibility_gap}
\end{equation}
is attained, and
\begin{equation}
 d_r^{\mathrm{feas}}=0
 \quad\Longleftrightarrow\quad
 \boldsymbol\mu\in\Omega_r.
 \label{eq:app_exact_realizability}
\end{equation}
Thus, if $d_r^{\mathrm{feas}}>0$, no admissible aggregation rule with the same
finite-window support and constraints can exactly realize
$\boldsymbol\mu$ in that window.  P2 should therefore be interpreted as
calibrating within the realizable image rather than synthesizing unavailable
target support.\hfill$\square$

\paragraph{Reference stability of P2.}
Fix one active window, including $M_r$, $\mathbf Q_r$, $V_r$,
$\overline{\boldsymbol\tau}_r$, $\mathcal C_r$, and all regularization
coefficients.  Let $\boldsymbol\alpha^\star(\boldsymbol\beta)$ denote the
unique P2 minimizer when the corrected reference is
$\boldsymbol\beta$.  Then, for any two reference vectors
$\boldsymbol\beta$ and $\boldsymbol\beta'$ on the same feasible support,
\begin{equation}
 \left\|
 \boldsymbol\alpha^\star(\boldsymbol\beta)
 -
 \boldsymbol\alpha^\star(\boldsymbol\beta')
 \right\|_2
 \le
 \left\|
 \boldsymbol\beta-\boldsymbol\beta'
 \right\|_2.
 \label{eq:app_reference_nonexpansive}
\end{equation}
To see this, write the P2 objective as
\begin{equation}
 G_r(\boldsymbol\alpha)
 +\frac{\lambda_\beta}{2}
 \left\|
 \boldsymbol\alpha-\boldsymbol\beta
 \right\|_2^2,
 \label{eq:app_p2_split}
\end{equation}
where $G_r$ contains the debt, group, variance, and staleness terms and is
convex under A4.  Let
$\mathbf x=\boldsymbol\alpha^\star(\boldsymbol\beta)$ and
$\mathbf y=\boldsymbol\alpha^\star(\boldsymbol\beta')$.  Adding the two
variational inequalities on the common convex set $\mathcal C_r$ gives
\begin{align}
 &\left\langle
 \nabla G_r(\mathbf x)-\nabla G_r(\mathbf y),
 \mathbf x-\mathbf y
 \right\rangle
 +\lambda_\beta\|\mathbf x-\mathbf y\|_2^2
 \nonumber\\
 &\qquad\le
 \lambda_\beta
 \left\langle
 \boldsymbol\beta-\boldsymbol\beta',
 \mathbf x-\mathbf y
 \right\rangle.
 \label{eq:app_reference_vi}
\end{align}
The first term on the left is nonnegative by monotonicity of the gradient of a
convex function.  Cauchy--Schwarz therefore yields
\eqref{eq:app_reference_nonexpansive}.
\qed

This property concerns only the P2 map
when the feasible set and all other window states are fixed.  It is not a
robustness guarantee to changes in the underlying mobile opportunity or
selection process.

\subsection{Proof-Obligation Summary}
\label{app:proof_summary}

The completed arguments establish the exact scope of the theory used in
the main text:

\begin{enumerate}
    \item Proposition~1 proves that unequal deployment and arrival laws define
    different risk functionals in general.
    \item Proposition~2 proves consistency of the ideal normalized sequential
    opportunity--observation--usable correction under A1--A3 for compatible
    target moments or target-group mass.  It does not prove optimization of
    deployment-target risk.
    \item Proposition~3 proves that, under A1, A4, A5, and A6b, the implemented
    clipped plug-in is a Lipschitz perturbation of the same-window ideal
    reference up to the measurable stabilization error
    $\varepsilon_{\mathrm{stab},r}$, and that P2 does not amplify that error on
    a fixed $\mathcal{A}_r$.  A6a is the additional clause that drives
    $\varepsilon_{\mathrm{stab},r}$ to zero.
    \item Theorem~1 proves existence and uniqueness of the solution to (P2)
    under A4.
    \item The debt recursion gives a deterministic long-term accounting bound.
    Vanishing structural mismatch follows whenever debt grows sublinearly in
    active step mass.
\end{enumerate}

None of these results asserts identification outside target support, robustness
to arbitrary hidden confounding, exact unbiasedness under active clipping, or
universal predictive superiority.


\section{Complete Proofs for Aggregate Coverage Certification}
\label{app:agg_coverage_proofs}

This section records the complete proofs for
Section~\mbox {III-D}--\mbox {VI-F}.
Existing reconstruction proofs are unchanged.  The coverage-certificate
numbering is independent of Proposition~1 and
Theorem~1.  All arguments below are complete.

\begin{algorithm}[t]
\caption{Aggregate coverage certificate (diagnostic Gate)}
\label{alg:agg_certificate}
\begin{algorithmic}[1]
\REQUIRE legal $F_{t-}$; public $a_t$; block snapshot $f_b$; same-block scores;
previous CERTIFIED miscoverage flag; $\alpha=0.05$; $\tau_{\mathrm{cov}}=0.05$; $n_{\min}=19$
\ENSURE $\mathrm{state}\in\{\mathrm{WARMUP},\mathrm{CERTIFIED},\mathrm{RESET}\}$,
block\_id, predictor\_snapshot\_id, $\widehat C_t$, $n$, $q_t$ if CERTIFIED,
$C^{\mathrm{LCB}}_t$, $\mathrm{Gate}_t$
\IF{the previous completed CERTIFIED window satisfied $C^{\mathrm{LCB}}>C$}
\STATE $\mathrm{state}\leftarrow\mathrm{RESET}$, $\mathrm{Gate}_t\leftarrow\mathrm{OFF}$
\STATE Fit $f_{b+1}$ on completed past data only; freeze it; clear the buffer
\STATE $\mathrm{state}\leftarrow\mathrm{WARMUP}$
\ELSE
\STATE Use the frozen snapshot $f_b$ of the current block
\IF{$n<19$ or $k_n>n$}
\STATE $\mathrm{state}\leftarrow\mathrm{WARMUP}$
\STATE $C^{\mathrm{LCB}}_t\leftarrow 0$, $\mathrm{Gate}_t\leftarrow\mathrm{OFF}$
\ELSE
\STATE $\mathrm{state}\leftarrow\mathrm{CERTIFIED}$ \COMMENT{scope; does not prove Assumption~1}
\STATE $\widehat C_t\leftarrow f_b(F_{t-})$
\STATE $q_t\leftarrow S_{(k_n)}^{\mathrm{cal}}$ \COMMENT{no clipping of $k_n$ to $n$}
\STATE $C^{\mathrm{LCB}}_t\leftarrow\max\{0,\widehat C_t-q_t\}$
\STATE $\mathrm{Gate}_t\leftarrow \mathbf{1}\{C^{\mathrm{LCB}}_t\ge\tau_{\mathrm{cov}}\}$
\ENDIF
\ENDIF
\end{algorithmic}
\end{algorithm}

\begin{proposition}[Positive pairwise LCB barrier]
\label{prop:pairwise-lcb-barrier}
Let $X_t\ge 0$ denote the realized opportunity of one target pair in window
$t$, and let $L_t=L(F_{t-})$ be a deterministic pre-window lower bound.  If
$P(X_t=0\mid F_{t-})>\alpha$, then no strictly positive $L_t>0$ can satisfy
$P(L_t\le X_t\mid F_{t-})\ge 1-\alpha$.
\end{proposition}

\begin{corollary}[Availability-induced pairwise barrier]
\label{cor:availability-barrier}
For a target pair $i=(k,s)$ let the designated stressed entity be $e(i)=k$
(the pair's client).  Let $A_{e(i),t}\in\{0,1\}$ be that entity's availability
in window $t$.  Assume $X_{i,t}>0\Rightarrow A_{e(i),t}=1$.  Under frozen
availability stress, $P(A_{e(i),t}=1\mid F_{t-},a)\le a$, hence
$P(X_{i,t}>0\mid F_{t-},a)\le a$.  If $a<1-\alpha$, no strictly positive
deterministic $1-\alpha$ pairwise LCB is guaranteed for that pair under
Proposition~\ref{prop:pairwise-lcb-barrier}.  The claim applies only to pairs
obeying the designated-entity necessary condition.
\end{corollary}

\begin{corollary}[False-safe Gate bound]
\label{cor:false-safe-gate}
Under the conditions of Theorem~2,
\begin{equation}
 P\bigl(\mathrm{Gate}_t=1\text{ and }C_t<\tau_{\mathrm{cov}}\mid g\bigr)
 \le\alpha.
 \label{eq:agg_cor2}
\end{equation}
\end{corollary}

\begin{proposition}[Arbitrary-drift applicability boundary]
\label{prop:drift-boundary}
Suppose two data-generating worlds share the same legal $F_{t-}$ but have
$C_t=c_{\mathrm{high}}$ and $C_t=c_{\mathrm{low}}$ with
$c_{\mathrm{low}}<\tau_{\mathrm{cov}}<c_{\mathrm{high}}$.  Any deterministic
pre-window certificate outputs the same $C^{\mathrm{LCB}}_t$ in both worlds.
If that output is at least $\tau_{\mathrm{cov}}$, the certificate is
false-safe in the low world; if it is below $\tau_{\mathrm{cov}}$, it cannot
certify the high world.  Without an assumption linking the current window to
observable pre-window information, no method can simultaneously provide
nontrivial per-window certification and universal protection against
arbitrary unannounced regime shifts.
\end{proposition}

\subsection{Proof of Proposition~\ref{prop:pairwise-lcb-barrier}}
\label{app:proof_pairwise_lcb}

Fix a version of the regular conditional probability given $F_{t-}$.  On the
event $\{L_t>0\}$, the inclusion $\{L_t\le X_t\}\subseteq\{X_t>0\}$ holds
pathwise: if $X_t=0$, then $X_t<L_t$, so $L_t\le X_t$ fails.  Therefore
\begin{align*}
P(L_t\le X_t\mid F_{t-})
&\le
P(X_t>0\mid F_{t-}) \\
&=
1-P(X_t=0\mid F_{t-}) \\
&<
1-\alpha
\end{align*}
on the set where $P(X_t=0\mid F_{t-})>\alpha$ and $L_t>0$.  Hence no strictly
positive deterministic $L(F_{t-})$ can meet the $1-\alpha$ conditional coverage
requirement on such $F_{t-}$.  Any deterministic lower bound that does meet
the requirement must satisfy $L_t=0$ on that $F_{t-}$, or the coverage claim
is false.  This is a structural non-vacuity barrier, not a statement that
pairwise diagnostics are useless.

\subsection{Proof of Corollary~\ref{cor:availability-barrier}}
\label{app:proof_availability_barrier}

Let $i=(k,s)$ and $e(i)=k$ be the designated stressed entity of the pair.
Assume $X_{i,t}>0\Rightarrow A_{e(i),t}=1$.  Then
$\{X_{i,t}>0\}\subseteq\{A_{e(i),t}=1\}$, so
\[
 P(X_{i,t}>0\mid F_{t-},a)
 \le
 P(A_{e(i),t}=1\mid F_{t-},a).
\]
Frozen entity thinning gives $P(A_{e(i),t}=1\mid F_{t-},a)\le a$, hence
$P(X_{i,t}>0\mid F_{t-},a)\le a$ and $P(X_{i,t}=0\mid F_{t-},a)\ge 1-a$.
If $a<1-\alpha$, then $P(X_{i,t}=0\mid F_{t-},a)>\alpha$, and
Proposition~\ref{prop:pairwise-lcb-barrier} applies to this pair.

The argument is scoped to pairs obeying the designated-entity necessary
condition.  It does not claim that unavailability of some unspecified member
of an arbitrary entity group yields $P(\mathrm{active})\le a$.

\subsection{$n_{\min}$ derivation}
\label{app:nmin}

A finite calibration quantile exists without clipping iff
$k_n=\lceil(n+1)(1-\alpha)\rceil\le n$.  If $(n+1)(1-\alpha)>n$, the ceiling is
at least $n+1$, so no finite order statistic of the calibration sample can
serve as $q_t$.  The inequality rearranges to $n+1<1/\alpha$, or
$n<1/\alpha-1$.  Thus a finite $q_t$ requires
$n\ge\lceil 1/\alpha\rceil-1$.  For $\alpha=0.05$, $n_{\min}=19$.  Clipping
$k_n$ to $n$ is exactly the $k_n=n+1$ / $q=+\infty$ case disguised as a finite
quantile; it does not restore $1-\alpha$ coverage for $n<n_{\min}$.  In
particular $n=8$ is not sufficient for a nominal $95\%$ finite-sample split
conformal guarantee.

\subsection{Proof of Theorem~2}
\label{app:proof_agg_lcb}

\paragraph{Auxiliary continuous tie-break and randomized rank.}
Let $S_{j_1},\ldots,S_{j_n}$ be the calibration scores and $S_t$ the current
score.  Assumption~1 says these $n+1$ random
variables are exchangeable given $g$.  Draw auxiliary variables
$U_{j_1},\ldots,U_{j_n},U_t$ i.i.d.\ Continuous$(0,1)$, independent of the
scores, and order the pairs $(S_j,U_j)$ lexicographically.  The $n+1$ pairs
are exchangeable given $g$ and almost surely distinct after tie-break, so the
randomized rank $R_t^\star$ of $(S_t,U_t)$ is uniform on $\{1,\ldots,n+1\}$:
$P(R_t^\star=r\mid g)=1/(n+1)$.  With
$k_n=\lceil(n+1)(1-\alpha)\rceil$,
\[
 P(R_t^\star\le k_n\mid g)=\frac{k_n}{n+1}\ge 1-\alpha.
\]
CERTIFIED enforces $k_n\le n$, so $q_t=S_{(k_n)}^{\mathrm{cal}}$ exists as a
calibration order statistic.

\paragraph{Event inclusion.}
We claim $\{R_t^\star\le k_n\}\subseteq\{S_t\le S_{(k_n)}^{\mathrm{cal}}\}$.
If $S_t>S_{(k_n)}^{\mathrm{cal}}$, then at least $k_n$ calibration scores are
strictly smaller than $S_t$, so the test pair ranks at least $k_n+1$ for every
$U_t$, contradicting $R_t^\star\le k_n$.  Hence the inclusion holds, and
\[
 P(S_t\le q_t\mid g)
 \ge
 P(R_t^\star\le k_n\mid g)
 \ge 1-\alpha.
\]
The opposite (invalid) direction is not used: if $A\subset B$ and
$P(B)\ge 1-\alpha$, one cannot conclude $P(A)\ge 1-\alpha$.

\paragraph{Score translation.}
By definition $S_t=\widehat C_t-C_t$, so
$\{S_t\le q_t\}=\{\widehat C_t-q_t\le C_t\}$.

\paragraph{Clipping at zero.}
Set $C^{\mathrm{LCB}}_t=\max\{0,\widehat C_t-q_t\}$.  Always $C_t\ge 0$.
If $\widehat C_t-q_t\le 0$, then $C^{\mathrm{LCB}}_t=0\le C_t$ holds pathwise,
whether or not $S_t\le q_t$.  If $\widehat C_t-q_t>0$, then
$C^{\mathrm{LCB}}_t=\widehat C_t-q_t$, and $\{S_t\le q_t\}$ is exactly
$\{C^{\mathrm{LCB}}_t\le C_t\}$.  In all cases
$\{S_t\le q_t\}\subseteq\{C^{\mathrm{LCB}}_t\le C_t\}$.  Clipping cannot
invalidate a valid lower bound.  Therefore
\[
 P\bigl(C^{\mathrm{LCB}}_t\le C_t\mid g\bigr)
 \ge
 P(S_t\le q_t\mid g)
 \ge 1-\alpha.
\]

\paragraph{Frozen predictor snapshot.}
$\widehat C_t=f_b(F_{t-})$ uses the frozen snapshot of block $b$.  Refitting
the scoring rule on every window with new outcomes would leave the ordinary
finite-sample exchangeable rank argument.  Current $C_t$ may enter the
same-block calibration buffer only after window $t$ completes.  The current
window is never used to form its own $q_t$.  Operational CERTIFIED does not
prove Assumption~1; the theorem remains
conditional on that assumption.

This completes the proof.

\subsection{Proof of Corollary~\ref{cor:false-safe-gate}}
\label{app:proof_false_safe}

On the event $\{\mathrm{Gate}_t=1,\,C_t<\tau_{\mathrm{cov}}\}$ one has
$C^{\mathrm{LCB}}_t\ge\tau_{\mathrm{cov}}>C_t$, hence
$C^{\mathrm{LCB}}_t>C_t$.  Therefore
\[
 \{\mathrm{Gate}_t=1,\,C_t<\tau_{\mathrm{cov}}\}
 \subseteq
 \{C^{\mathrm{LCB}}_t>C_t\}.
\]
Theorem~2 gives $P(C^{\mathrm{LCB}}_t\le C_t\mid g)\ge 1-\alpha$,
so $P(C^{\mathrm{LCB}}_t>C_t\mid g)\le\alpha$.  The false-safe probability is
at most $\alpha$.  In WARMUP and RESET, $\mathrm{Gate}_t$ is OFF by definition,
so the false-safe event is empty on those windows.

\subsection{Proof of Proposition~\ref{prop:drift-boundary}}
\label{app:proof_drift_boundary}

Let $\mathcal{A}$ be any map from $F_{t-}$ to a real number
$C^{\mathrm{LCB}}_t$, possibly using extra bits independent of the current
world.  Because Worlds A and B share $F_{t-}$, the law of
$C^{\mathrm{LCB}}_t$ is identical in both worlds.  Write $L$ for that common
output.  Exactly one of $L\ge\tau_{\mathrm{cov}}$ or $L<\tau_{\mathrm{cov}}$
holds for each realization.  If $L\ge\tau_{\mathrm{cov}}$, then in World B one
has $L\ge\tau_{\mathrm{cov}}>c_{\mathrm{low}}=C_t$, so Gate ON is false-safe.
If $L<\tau_{\mathrm{cov}}$, then in World A the Gate is OFF even though
$C_t=c_{\mathrm{high}}>\tau_{\mathrm{cov}}$, so the window is not certified.
No choice of $L=L(F_{t-})$ avoids both failures.  This is an applicability
boundary, not a defect of a particular predictor.  A detector that uses only
$F_{t-}$ cannot protect the first unannounced shift window if that window is
indistinguishable from the certified regime in $F_{t-}$.  After a detected
change, RESET then WARMUP (Gate OFF) until $n\ge 19$ in the new regime is the
correct operational response.  The detector does not restore
Theorem~2 on the undetected first-shift window.

The optional bounded-TV robustness extension is omitted: it would require a
legally frozen pre-test $\epsilon_g$ with $\alpha'+\epsilon_g\le\alpha$, not
estimated post hoc from the test sequence, plus a complete proof.  Those
conditions are not met, so the core paper relies only on
Theorem~2, Corollary~\ref{cor:false-safe-gate}, and
Proposition~\ref{prop:drift-boundary}.


\section{Supplementary Risk-Semantic Evidence}
\label{app:risk_semantic_evidence}

This appendix reports supplementary statistics that are useful for interpreting
risk interchangeability but are not required in the main text.  No additional realization,
retraining, or post-hoc model selection is introduced here.

\subsection{Balanced No-Harm Control}
\label{app:e1_balanced}

Table~\ref{tab:app_e1_balanced} reports all five paired methods under the balanced
Balanced control.  The selected TimeAlign reference has mean
$\mathrm{RMSE}_\mu=7.7108$, while RAVEN-MCS obtains $7.7154$.  RAVEN's
pre-specified mean relative degradation is $0.0599\%$; the one-sided 95\%
upper bound is $0.134\%$, below the pre-specified $3\%$ no-harm margin.  The
balanced experiment is therefore a sanity check, not evidence of statistical
superiority.

\begin{table}[!t]
\centering
\caption{Balanced control over the five pre-specified paired seeds.}
\label{tab:app_e1_balanced}
\footnotesize
\setlength{\tabcolsep}{5pt}
\begin{tabular*}{\columnwidth}{@{\extracolsep{\fill}}lcc}
\toprule
Method & $\mathrm{RMSE}_\mu$ & Seeds\\
\midrule
FedAvg-Window & $7.7220\pm0.0289$ & 5 \\
FedAsync-Window & $7.7222\pm0.0287$ & 5 \\
TimeAlign-Agg & $7.7108\pm0.0272$ & 5 \\
TwoStage-H\'ajek & $7.7262\pm0.0291$ & 5 \\
RAVEN-MCS & $7.7154\pm0.0337$ & 5 \\
\bottomrule
\end{tabular*}
\end{table}

\subsection{Complete Risk Non-Interchangeability Statistics}
\label{app:e2_full_stats}

Table~\ref{tab:app_e2_daligned} gives the complete method-level risk-interchangeability summary for
the pre-specified difference-of-differences
$D_{\mathrm{aligned}}=\Delta\mathrm{RMSE}_{\mu}
-\Delta\mathrm{RMSE}_{\rho}$.  The negative mean is stable across all five
methods, and 17--18 of 20 paired seeds are negative.  This supports the
risk-semantic premise that arrival-weighted and deployment-target-weighted
views react differently to the same structured selection shift.  It does not
establish method-level superiority.

\begin{table}[!t]
\centering
\caption{Complete-aligned risk non-interchangeability statistics.
CI denotes the 95\% interval; $p$ is the paired Wilcoxon test of
$D_{\mathrm{aligned}}$ against zero.}
\label{tab:app_e2_daligned}
\scriptsize
\setlength{\tabcolsep}{2.5pt}
\begin{tabular*}{\columnwidth}{@{\extracolsep{\fill}}lcccc}
\toprule
Method & $D_{\mathrm{aligned}}$ & 95\% CI & Neg. seeds & $p$\\
\midrule
FedAsync-Window & $-0.02115$ & $[-0.03158,-0.01073]$ & 18/20 & $4.83e-04$ \\
FedAvg-Window & $-0.02115$ & $[-0.03158,-0.01071]$ & 18/20 & $4.83e-04$ \\
TimeAlign-Agg & $-0.02137$ & $[-0.03185,-0.01090]$ & 17/20 & $4.83e-04$ \\
RAVEN-MCS & $-0.02088$ & $[-0.03132,-0.01043]$ & 17/20 & $5.86e-04$ \\
TwoStage-H\'ajek & $-0.02120$ & $[-0.03169,-0.01072]$ & 17/20 & $5.86e-04$ \\
\bottomrule
\end{tabular*}
\end{table}


\section{Local-H\'ajek Degeneracy Analysis}
\label{app:local_hajek}

The equality between Local-Hajek and FedAvg in the formal SensorScope/U-Air
blocks is protocol-specific and should not be interpreted as a general
equivalence.  The provenance audit covers
9666 effective client-windows: 6070 on SensorScope and
3596 on U-Air.  In every audited effective client-window, the realized
within-client correction-weight range is 0.

\begin{table}[!t]
\centering
\caption{Local-Hajek provenance audit.}
\label{tab:app_local_hajek}
\footnotesize
\setlength{\tabcolsep}{5pt}
\begin{tabular*}{\columnwidth}{@{\extracolsep{\fill}}lcc}
\toprule
Dataset & Effective client-windows & Weight range\\
\midrule
SensorScope & 6070 & $0$\\
U-Air & 3596 & $0$\\
\midrule
Total & 9666 & $0$\\
\bottomrule
\end{tabular*}
\end{table}

The collapse follows directly from the normalized local objective.  Let
$\mathcal I_{k,r}^{+}$ be the records retained by an effective client-window
and let $N_{k,r}^{+}=|\mathcal I_{k,r}^{+}|$.  If the realized first-stage
weight is constant over that set, i.e.,
$a_{k,r,i}=a_{k,r}>0$ for all $i\in\mathcal I_{k,r}^{+}$, then
\begin{equation}
 \overline a_{k,r,i}
 =
 \frac{a_{k,r}}
 {N_{k,r}^{+}a_{k,r}}
 =
 \frac{1}{N_{k,r}^{+}}.
 \label{eq:app_local_hajek_uniform}
\end{equation}
Consequently,
\begin{equation}
 \widehat F^H_{k,r}(\theta)
 =
 \frac{1}{N_{k,r}^{+}}
 \sum_{i\in\mathcal I_{k,r}^{+}}
 \frac{1}{2}
 \left(
 f_\theta(i;C_i^{\mathrm{dep}})-Z_{k,r,i}
 \right)^2,
 \label{eq:app_local_hajek_collapse}
\end{equation}
which is exactly the uniform local objective used by the matched FedAvg
backbone in these realized windows.
\qed

This derivation explains the observed
numerical equality without weakening the general role of H\'ajek normalization
when within-window correction weights are heterogeneous.


\section{P2 Feasibility Diagnostics}
\label{app:p2_diagnostics}

Theorem~1 concerns the mathematical program.  Separately, the
experimental audit reconstructs whether logged allocations satisfy the P2
constraints whenever the necessary state is available.  Table~\ref{tab:app_p2_coverage}
summarizes the dataset-level certificate coverage.

\begin{table}[!t]
\centering
\caption{P2 feasibility-certificate coverage in the formal protocol blocks.}
\label{tab:app_p2_coverage}
\footnotesize
\setlength{\tabcolsep}{4pt}
\begin{tabular*}{\columnwidth}{@{\extracolsep{\fill}}lcccc}
\toprule
Dataset & Runs & Mean windows & Range & Status\\
\midrule
SensorScope & 180 & 93.85 & 88--98 & Certified \\
U-Air & 180 & 72.61 & 63--79 & Certified \\
\bottomrule
\end{tabular*}
\end{table}

The solver log contains 778 numerical-inaccuracy warnings across
85 affected formal runs.  The reconstruction audit covers
6438 active P2 windows and certifies feasibility; the maximum
reconstructed constraint residual is
$2.686\times10^{-8}<10^{-7}$.  Feasibility is therefore certified;
exact post-hoc optimality of every logged solve is not claimed.
On later reconstructable logs, CLARABEL and SCS typically match to
$10^{-8}$ in $\ell_1$; the worst window differs by $0.0056$ in
allocation, which does not move the main-comparison or worst-group
headlines.  A pre-registered window-level
$10^{-4}/10^{-3}$ divergence rule is recorded as
\texttt{G5\_SOLVER\_DIVERGE}; the stored P2 outputs were not re-solved under a different solver.

\begin{table}[!t]
\centering
\caption{Warning-affected P2 audit and claim boundary.}
\label{tab:app_p2_warning_audit}
\footnotesize
\setlength{\tabcolsep}{5pt}
\begin{tabular*}{\columnwidth}{@{\extracolsep{\fill}}lc}
\toprule
Quantity & Frozen result\\
\midrule
Numerical-inaccuracy warnings & 778\\
Affected formal runs & 85\\
Audited active windows & 6438\\
Maximum constraint residual & $2.686\times10^{-8}$\\
Feasibility & Certified\\
Exact post-hoc optimality & Not claimed\\
\bottomrule
\end{tabular*}
\end{table}


\FloatBarrier

\section{Additional Existing Experimental Context}
\label{app:additional_context}

This appendix collects additional context that is useful for interpretation but
would crowd the main RQ-focused evaluation.  Reference-only methods are marked
by $^\dagger$ and are not included in the main inferential family.

The sealed same-backbone protocol family below uses the Design-off setting on U-Air;
the main-text table reports both settings on both blocks, and the
sensor-scope and U-Air rows of Table~III are the Design-on and Design-off
settings respectively.

\subsection{Windowed adaptations of FedAU, ObsUse, and TimeAlign}
\label{app:adapted_baselines}
FedAU-Window, ObsUse-Window, and TimeAlign-Agg are window-close
adaptations, not official reproductions.
FedAU~\cite{wang2024lightweight} targets the original federated averaging
objective under unknown participation; these rules reweight usable clients
and do not recover that objective.
After window close, each EventTrace record of client \(k\) increments
\(N^{\mathrm{att}}_k\), and \(N^{\mathrm{use}}_k\) if it is a registered
usable update.  Window \(r\) uses earlier windows
(\(\varepsilon=10^{-6}\); unseen rate \(1\)).
Let \(n_{k,r}\) be the observed-unit count (empty buffer encoded as one),
\(\widehat p^{\mathrm{obs}}_{k,r}\) the mean frozen observation propensity
on the current risk set clipped below at \(0.05\) (or \(1\) if empty),
\(S_k\) the covered slots on \(\mathcal{A}_r\), and
\(\mathrm{overlap}(s)=\#\{j\in\mathcal{A}_r:s\in S_j\}\).
\begin{align}
 \widehat p^{\mathrm{part}}_{k,r}
 &=
 (N^{\mathrm{use}}_k+\varepsilon)/(N^{\mathrm{att}}_k+\varepsilon),
 \label{eq:fedau_part_rate}\\
 \alpha^{\mathrm{FedAU}}_{k,r}
 &\propto n_{k,r}/\max\{\widehat p^{\mathrm{part}}_{k,r},\varepsilon\},
 \label{eq:fedau_window_alpha}\\
 \alpha^{\mathrm{ObsUse}}_{k,r}
 &\propto n_{k,r}/(\widehat p^{\mathrm{obs}}_{k,r}\widehat p^{\mathrm{part}}_{k,r}),
 \label{eq:obsuse_window_alpha}\\
 \alpha^{\mathrm{TimeAlign}}_{k,r}
 &\propto \sum_{s\in S_k}1/\mathrm{overlap}(s).
 \label{eq:timealign_window_alpha}
\end{align}
FedAU-Window uses \eqref{eq:fedau_window_alpha} only.
ObsUse-Window uses \eqref{eq:obsuse_window_alpha} and omits Design, (P2), and debt.
TimeAlign-Agg (\texttt{flamf\_timealign\_adapted}) uses
\eqref{eq:timealign_window_alpha} following~\cite{wang2026federated},
without \(n_{k,r}\) or Design, and is not official FLAMF.
Paired methods share the Common-NDMF backbone, step size, and selection
trace (seeds \(30001\)--\(30020\)).  Existing tables already report
these adaptations.

\begin{table*}[t]
\centering
\caption{Nearest-neighbor comparison.  RAVEN's distinction is not IPW
itself, but a public deployment target, sequential R/O/U legal
histories, sensing-side correction, and finite-window calibration.
FedIPW~\cite{morishita2026who} is the nearest two-stage FL IPW
(enrollment then participation); it is not an MCS baseline.
FLOSS~\cite{goetze2025floss} is client-level opt-out/straggler IPW.
FedAU-Window and ObsUse-Window used in Section~VII are
windowed implementations, not official reproductions.
This table is not mixed into Table~III.}
\label{tab:nearest_neighbor}
\setlength{\tabcolsep}{3.2pt}
{\footnotesize
\begin{tabular}{@{}p{2.30cm}p{2.10cm}p{2.05cm}p{1.85cm}p{2.30cm}@{}}
\toprule
Family & Target & Selection & Sensing-side & Finite-support cal.\\
\midrule
Covariate-shift FL & target/test domain & sample/client & no & no\\
FedAU / participation debiasing & original FL objective & client participation & no & no\\
FedIPW & pop.\ mean update & enroll then participate & no & no\\
FLOSS & client roster & opt-out / straggler & no & no\\
Correlated participation & long-run availability & client/update & no & limited\\
Agnostic / targeted FL & client mixture & client distribution & no & not the issue\\
Semi-async fairness/queue & participation fairness & client/update & no & queue \(\neq\) target debt\\
\textbf{RAVEN-MCS} & public deployment \(\mu\) & R/O/U & yes & yes\\
\bottomrule
\end{tabular}}
\end{table*}

\begin{table}[t]
\centering
\caption{Validation-only under-served-block RMSE of the Table~III
RAVEN configuration relative
to FedAU-Window and ObsUse-Window (ten paired seeds; negative means
the RAVEN row is better).  The reported settings are Design on for SensorScope and Design off for U-Air.
A win requires both columns $\le-0.5\%$.  Neither
dataset meets that rule; both satisfy no-harm ($\le+0.5\%$).  The test
look was not opened.  Not mixed into Table~III.}
\label{tab:underserved_block}
\setlength{\tabcolsep}{3.5pt}
\begin{tabular}{lcccc}
\toprule
Dataset & $g^{\star}$ (UTC) & vs FedAU & vs ObsUse & Win / no-harm\\
\midrule
SensorScope & 2 (12--18) & $+0.05\%$ & $+0.15\%$ & No / Yes\\
U-Air & 3 (18--24) & $-0.44\%$ & $-0.27\%$ & No / Yes\\
\bottomrule
\end{tabular}
\end{table}

\begin{figure}[t]
\centering
\includegraphics[width=\columnwidth]{fig3_e3_structural_mismatch.pdf}
\caption{Sealed same-backbone structural alignment on SensorScope and
U-Air (mean $\pm$ std over 20 paired seeds; lower is better).
(a)(b)~SensorScope $\Delta_{\mathrm{group}}$ and $\Delta_{c\text{-}s}$;
(c)(d)~U-Air $\Delta_{\mathrm{group}}$ and $\Delta_{c\text{-}s}$.
SensorScope RAVEN-MCS is RAVEN-Design.  U-Air RAVEN-MCS in this
figure uses Design on, not the Design-off rows of
Table~III.}
\label{fig:e3_structural}
\end{figure}

\begin{figure}[t]
\centering
\includegraphics[width=\columnwidth]{fig4_e3_prediction_risk.pdf}
\caption{Sealed same-backbone deployment-target prediction risk
(mean $\pm$ std over 20 paired seeds; lower is better).
(a)(b)~SensorScope $\mathrm{RMSE}_{\mu}$ and Tail RMSE;
(c)(d)~U-Air $\mathrm{RMSE}_{\mu}$ and Tail RMSE.
U-Air RAVEN-MCS here uses Design on.  The Design-off U-Air row
($\mathrm{RMSE}_{\mu}=1.1131$) is in
Table~III.}
\label{fig:e3_prediction}
\end{figure}

\begin{figure}[t]
\centering
\includegraphics[width=\columnwidth]{fig5_e4_component_ablation.pdf}
\caption{One-factor component removal.  Positive bars mean that removing the
component worsens the metric; $^\ast$ marks Holm-adjusted $p<0.05$.
(a)(b)~SensorScope structural and prediction endpoints;
(c)(d)~U-Air structural and prediction endpoints.
Values with $|\delta|<0.05\%$ and no significance mark are left unlabeled.
Sign reversals show endpoint- and dataset-dependent roles rather than
uniformly additive benefits.  The corresponding numeric table remains
Table~IV in the main paper.}
\label{fig:e4_ablation}
\end{figure}

\begin{figure}[t]
\raggedright
\includegraphics[width=\columnwidth]{fig6_e4_uair_excl_design.pdf}
\caption{U-Air component ablation excluding Design, magnified from
Fig.~\ref{fig:e4_ablation}.  (a)~Structural mismatch;
(b)~prediction risk.  Positive means removing the component worsens
the metric; $^\ast$ marks Holm-adjusted $p<0.05$.}
\label{fig:uair_excl_design}
\end{figure}

\subsection{Full Method Context on the Comparable Datasets}
\label{app:e3_full_context}

\begin{table}[!htbp]
\centering
\caption{SensorScope full method context (mean $\pm$ std over 20 seeds).
Lower is better.  $^\dagger$ denotes reference-only context.
RAVEN-Design is the Table~III method.}
\label{tab:app_e3_sensorscope}
\scriptsize
\setlength{\tabcolsep}{3pt}
\begin{tabular*}{\columnwidth}{@{\extracolsep{\fill}}lccc}
\toprule
Method & $\mathrm{RMSE}_{\mu}$ & $\Delta_{\mathrm{group}}$ & $\Delta_{c\text{-}s}$\\
\midrule
Central-All$^{\dagger}$ & $7.6004\pm0.0695$ & -- & -- \\
Central-Delivered$^{\dagger}$ & $7.9291\pm0.0759$ & -- & -- \\
FedAvg-Window & $7.9389\pm0.0459$ & $1.0558\pm0.0288$ & $1.2903\pm0.0306$ \\
FedAsync-Window & $7.9391\pm0.0457$ & $1.0558\pm0.0287$ & $1.2906\pm0.0307$ \\
TimeAlign-Agg & $7.9395\pm0.0455$ & $1.0561\pm0.0282$ & $1.2904\pm0.0304$ \\
Local-Hajek & $7.9389\pm0.0459$ & $1.0558\pm0.0288$ & $1.2903\pm0.0306$ \\
TwoStage-H\'ajek & $7.9309\pm0.0543$ & $1.0353\pm0.0347$ & $1.2910\pm0.0324$ \\
Inst-Cal & $7.9380\pm0.0473$ & $1.0141\pm0.0305$ & $1.2644\pm0.0312$ \\
Debt-Cal & $7.9399\pm0.0463$ & $1.0272\pm0.0275$ & $1.2710\pm0.0298$ \\
RAVEN-Design & $7.9380\pm0.0473$ & $1.0139\pm0.0305$ & $1.2644\pm0.0312$ \\
RAVEN-SimOracle$^{\dagger}$ & $7.9373\pm0.0455$ & $1.0160\pm0.0327$ & $1.2709\pm0.0301$ \\
\bottomrule
\end{tabular*}
\end{table}

\begin{table*}[t]
\centering
\caption{U-Air full method context (mean $\pm$ std over 20 seeds).
Lower is better.  $^\dagger$ denotes reference-only context.
Each row uses the active set its own definition licenses: the rows that read
$\pi^{\mathrm{tar}}$ restrict the set, the others keep every delivered client.
The Design-on row is not the Design-off setting reported in Table~III.}
\label{tab:app_e3_uair}
\scriptsize
\setlength{\tabcolsep}{3pt}
\begin{tabular*}{\textwidth}{@{\extracolsep{\fill}}lccc}
\toprule
Method & $\mathrm{RMSE}_{\mu}$ & $\Delta_{\mathrm{group}}$ & $\Delta_{c\text{-}s}$\\
\midrule
Central-All$^{\dagger}$ & $1.0556\pm0.0566$ & -- & -- \\
Central-Delivered$^{\dagger}$ & $1.1782\pm0.0981$ & -- & -- \\
FedAvg-Window & $1.1119\pm0.0603$ & $1.0446\pm0.0397$ & $1.3138\pm0.0337$ \\
FedAsync-Window & $1.1119\pm0.0599$ & $1.0451\pm0.0397$ & $1.3137\pm0.0336$ \\
TimeAlign-Agg & $1.1120\pm0.0597$ & $1.0448\pm0.0404$ & $1.3145\pm0.0340$ \\
Local-Hajek & $1.1923\pm0.0629$ & $1.1218\pm0.0417$ & $1.2621\pm0.0309$ \\
TwoStage-H\'ajek & $1.2202\pm0.0672$ & $1.1196\pm0.0445$ & $1.2406\pm0.0297$ \\
Inst-Cal & $1.1954\pm0.0635$ & $1.0924\pm0.0479$ & $1.2334\pm0.0320$ \\
Debt-Cal & $1.1144\pm0.0645$ & $1.0175\pm0.0481$ & $1.2948\pm0.0304$ \\
RAVEN-Design & $1.1954\pm0.0635$ & $1.0924\pm0.0479$ & $1.2334\pm0.0320$ \\
RAVEN-SimOracle$^{\dagger}$ & $1.1958\pm0.0630$ & $1.0922\pm0.0483$ & $1.2342\pm0.0325$ \\
\bottomrule
\end{tabular*}
\end{table*}

The central methods provide prediction-only reference context and therefore
have no distributed structural-mismatch entry.  RAVEN-SimOracle is a controlled
reference rather than a deployable baseline.  NSW-Traffic is intentionally not
added to this homogeneous method table because its reused cross-vintage path is not
scientifically comparable to the SensorScope/U-Air formal blocks.

\subsection{NSW-Traffic Four-Region Protocol Ranking}
\label{app:traffic_four_region}

Table~\ref{tab:traffic_four_region} ranks the same windowed family on
NSW-Traffic under a new public target map: four pre-specified coordinate
regions, not the UTC H1--H4 groups of Table~III.
Design is configured on for this block.
$R$, $O$, and $U$ remain Complete-aligned protocol layers.
This block is a new public target map on NSW-Traffic, not a third-city
MCS system; it is not mixed into Table~III;
and it is not comparable to the excluded cross-vintage Traffic path.
$\mathrm{WorstGroupRMSE}$ is the maximum error over the four regions.
RAVEN-Design versus FedAU-Window and ObsUse-Window stays within
the descriptive $+0.5\%$ worst-group mean-relative point-estimate
tolerance (paired relative $-0.57\%$ and $-0.54\%$).
FedAU-Window and ObsUse-Window remain smaller on $\Delta_{\mathrm{group}}$.
The block applies the named heuristic to a local map; it does not transfer
the SensorScope/U-Air support statistic.

\begin{table}[t]
\centering
\caption{NSW-Traffic ranking under a four-region public target map
(mean $\pm$ std, 10 seeds, $100$ windows).  Lower is better.
TwoStage-H\'ajek is lowest on $\mathrm{RMSE}_{\mu}$, WorstGroupRMSE, and
$\Delta_{c\text{-}s}$; FedAU-Window is lowest on
$\Delta_{\mathrm{group}}$.  RAVEN-Design stays within the
descriptive $+0.5\%$ worst-group tolerance versus FedAU-Window and
ObsUse-Window.}
\label{tab:traffic_four_region}
\scriptsize
\setlength{\tabcolsep}{2.5pt}
\begin{tabular}{@{}lcccc@{}}
\toprule
Method & $\mathrm{RMSE}_{\mu}$ & WorstGroupRMSE &
$\Delta_{\mathrm{group}}$ & $\Delta_{c\text{-}s}$\\
\midrule
FedAvg-Window & \shortstack{$1.0180$\\$\pm0.0089$} & \shortstack{$1.4245$\\$\pm0.0172$} & \shortstack{$0.1877$\\$\pm0.0263$} & \shortstack{$0.9597$\\$\pm0.0290$}\\
TwoStage-H\'ajek & \shortstack{$\mathbf{1.0095}$\\$\pm0.0052$} & \shortstack{$\mathbf{1.4025}$\\$\pm0.0128$} & \shortstack{$0.1908$\\$\pm0.0233$} & \shortstack{$\mathbf{0.8362}$\\$\pm0.0335$}\\
FedAU-Window & \shortstack{$1.0131$\\$\pm0.0078$} & \shortstack{$1.4110$\\$\pm0.0185$} & \shortstack{$\mathbf{0.0606}$\\$\pm0.0192$} & \shortstack{$0.9612$\\$\pm0.0333$}\\
ObsUse-Window & \shortstack{$1.0130$\\$\pm0.0078$} & \shortstack{$1.4106$\\$\pm0.0187$} & \shortstack{$0.0607$\\$\pm0.0192$} & \shortstack{$0.9602$\\$\pm0.0342$}\\
RAVEN-Design & \shortstack{$1.0099$\\$\pm0.0061$} & \shortstack{$1.4029$\\$\pm0.0156$} & \shortstack{$0.1048$\\$\pm0.0242$} & \shortstack{$0.8414$\\$\pm0.0304$}\\
\bottomrule
\end{tabular}

\begin{minipage}{\columnwidth}
\raggedright\scriptsize
$R$, $O$, and $U$ are Complete-aligned protocol layers.
WorstGroupRMSE is the maximum over four coordinate regions, not UTC
H1--H4.  RAVEN-Design is used (Design on).
Not mixed into Table~III; not a confirmation of
that table and not an end-to-end MCS ranking.
No Holm-significant win is claimed.
\end{minipage}
\end{table}

\subsection{Four-Block Per-Group RMSE}
\label{app:block_rmse}

Table~\ref{tab:app_block_rmse} reports the four UTC-block scores whose
maximum is $\mathrm{WorstGroupRMSE}$ in
Table~III.  The blocks are the public
deployment groups $00$--$06$, $06$--$12$, $12$--$18$, and $18$--$24$~UTC.
This table does not introduce a new ranking rule.  On both datasets the
worst block is $g=2$ ($12$--$18$~UTC).  The pre-registered under-served
block $g^{\star}$ coincides with that worst block on SensorScope and
does not on U-Air.

\begin{table*}[t]
\centering
\caption{Four-block $\mathrm{RMSE}_g$ (mean $\pm$ std over 20 paired
seeds).  Lower is better.  $\mathrm{WorstGroupRMSE}=\max_g\mathrm{RMSE}_g$.
Both Design settings are reported on both blocks.
Source: G2 stored per-seed block scores.}
\label{tab:app_block_rmse}
\footnotesize
\setlength{\tabcolsep}{4pt}
\begin{tabular*}{\textwidth}{@{\extracolsep{\fill}}llccccc}
\toprule
Dataset & Method & $g{=}0$ & $g{=}1$ & $g{=}2$ & $g{=}3$ & WG\\
\midrule
\multirow{6}{*}{SS}
& FedAvg-Window & $7.4912{\pm}0.0464$ & $7.5873{\pm}0.0444$ & $8.6716{\pm}0.0471$ & $8.2084{\pm}0.0460$ & $8.6716{\pm}0.0471$\\
& TwoStage-H\'ajek & $7.4831{\pm}0.0549$ & $7.5795{\pm}0.0523$ & $8.6633{\pm}0.0556$ & $8.2004{\pm}0.0549$ & $8.6633{\pm}0.0556$\\
& FedAU-Window & $7.4887{\pm}0.0512$ & $7.5848{\pm}0.0491$ & $8.6691{\pm}0.0522$ & $8.2059{\pm}0.0510$ & $8.6691{\pm}0.0522$\\
& ObsUse-Window & $7.4815{\pm}0.0517$ & $7.5780{\pm}0.0496$ & $8.6618{\pm}0.0528$ & $8.1988{\pm}0.0516$ & $8.6618{\pm}0.0528$\\
& RAVEN-Design & $7.4903{\pm}0.0479$ & $7.5864{\pm}0.0457$ & $8.6707{\pm}0.0486$ & $8.2075{\pm}0.0476$ & $8.6707{\pm}0.0486$\\
& RAVEN-NoDesign & $7.4907{\pm}0.0474$ & $7.5868{\pm}0.0452$ & $8.6711{\pm}0.0481$ & $8.2079{\pm}0.0470$ & $8.6711{\pm}0.0481$\\
\midrule
\multirow{6}{*}{UA}
& FedAvg-Window & $1.1015{\pm}0.0546$ & $1.1144{\pm}0.0578$ & $1.1690{\pm}0.0692$ & $1.0721{\pm}0.0601$ & $1.1690{\pm}0.0692$\\
& TwoStage-H\'ajek & $1.2009{\pm}0.0623$ & $1.2194{\pm}0.0649$ & $1.2907{\pm}0.0753$ & $1.1800{\pm}0.0670$ & $1.2907{\pm}0.0753$\\
& FedAU-Window & $1.1023{\pm}0.0545$ & $1.1153{\pm}0.0576$ & $1.1701{\pm}0.0689$ & $1.0729{\pm}0.0598$ & $1.1701{\pm}0.0689$\\
& ObsUse-Window & $1.1016{\pm}0.0544$ & $1.1145{\pm}0.0576$ & $1.1692{\pm}0.0689$ & $1.0721{\pm}0.0598$ & $1.1692{\pm}0.0689$\\
& RAVEN-Design & $1.1779{\pm}0.0583$ & $1.1953{\pm}0.0611$ & $1.2630{\pm}0.0716$ & $1.1552{\pm}0.0633$ & $1.2630{\pm}0.0716$\\
& RAVEN-NoDesign & $1.1025{\pm}0.0578$ & $1.1155{\pm}0.0611$ & $1.1703{\pm}0.0734$ & $1.0732{\pm}0.0636$ & $1.1703{\pm}0.0734$\\
\bottomrule
\end{tabular*}
\end{table*}

\subsection{Runtime Diagnostic from Existing Risk-Semantic Logs}
\label{app:runtime}

Table~\ref{tab:app_runtime} reports measured wall-clock runtime only as an
implementation diagnostic.  The ``all scenarios'' column aggregates six
risk-semantic scenarios and 20 seeds (120 runs per method); the Complete-aligned column
contains 20 runs per method.  These measurements are not a scalability claim.

\begin{table}[!htbp]
\centering
\caption{Existing risk-semantic runtime diagnostic in seconds (mean $\pm$ std).}
\label{tab:app_runtime}
\footnotesize
\setlength{\tabcolsep}{3.5pt}
\begin{tabular*}{\columnwidth}{@{\extracolsep{\fill}}lcc}
\toprule
Method & All scenarios & Complete-aligned\\
\midrule
FedAvg-Window & $71.426\pm5.105$ & $69.278\pm2.287$ \\
FedAsync-Window & $71.418\pm4.863$ & $70.186\pm3.531$ \\
TimeAlign-Agg & $71.244\pm5.225$ & $69.217\pm2.554$ \\
TwoStage-H\'ajek & $71.781\pm8.017$ & $69.829\pm3.670$ \\
RAVEN-MCS & $72.669\pm4.243$ & $71.070\pm1.292$ \\
\bottomrule
\end{tabular*}
\end{table}

Across these runtime rows (Table~\ref{tab:app_runtime}) there are no recorded
non-finite model parameters, non-finite losses, denominator-zero events, or
recorded P2 solver failures.  The diagnostic supports implementation sanity only and does not
establish asymptotic system efficiency.

\subsection{Provenance of the w/o-Design Endpoint Trade-off}
\label{app:design_provenance}

The Design ablation in Section~\mbox {VII-E} exhibits an endpoint
trade-off on U-Air.  A paired-window provenance audit covers
774 windows.  Opportunity, observation, usable-update, and
registration hashes are identical between full RAVEN and w/o-Design throughout
the audited set.  Differences first arise from the positive design-weight/P2
support induced after the Design ratio is removed.  This confirms that the
ablation effect is a weight-support consequence of the intended intervention,
not selection-trace leakage.

\section{Intermittent Protocol Opportunity (INT-R)}
\label{app:int_r}

The main comparison and ablation ranking in Table~III uses
Complete-aligned protocol opportunity that is always-on: every mapped
client receives an opportunity event in every window.  That setting does
not contain mobile clients, and it does not use T-Drive trajectories.
This appendix reports a robustness ranking on the same SensorScope and
U-Air station arrays, the same public four-block $\boldsymbol\mu$, and the
same Complete-aligned observation and usable-update layers, after replacing
always-on protocol $R$ by a client-window occupancy bit $A_{k,t}$.
If $A_{k,t}=0$, that client-window writes no opportunity event.  If
$A_{k,t}=1$, the existing Complete-aligned $O$/$U$ path is unchanged.
Occupancy is a two-state Markov process, independent across the $K=8$
logical clients, with mean active-fraction target $0.425$ taken from the
already-published T-Drive scalar in Section~\mbox {VII-B}.  The
same $(\alpha,\beta)$ pair is used on both datasets.  T-Drive GPS is not
overlaid on the stations.  The Design settings are not retuned from
these RMSE values: they are configurations, and both are reported, so
SensorScope and U-Air are shown with the settings recorded in Table~III.
The coverage Gate audited in Section~\mbox {VII-F} is not used.

A Goh--Barab\'asi burstiness target of $+0.38$ was also declared.  The
stability screen required to keep single-seed active fractions inside
$[0.35,0.50]$ selected a less sticky chain
($\alpha\approx0.775$, $\beta\approx0.837$) whose pooled burstiness is
about $0.20$.  This block therefore matches the published T-Drive
\emph{active fraction}, not T-Drive burstiness, and it is not a taxi
opportunity process.

Table~\ref{tab:int_r_occ} records the occupancy audit on seeds
$30001$, $30010$, and $30020$.  Mean active fraction on seed $30001$
lies in $[0.35,0.50]$ on both datasets.  Four-block
$\mathrm{TV}(\pi^{\mathrm{opp}},\boldsymbol\mu)$ remains above the
pre-registered $0.05$ threshold ($0.309$ on SensorScope and $0.225$ on
U-Air at seed $30001$), so intermittent occupancy does not erase the
mismatch that Table~III is about.

\begin{table}[!htbp]
\centering
\caption{INT-R occupancy audit on the shared EventTraces (eight logical
clients).  Burstiness is Goh--Barab\'asi on pooled inter-active gaps.
$\mathrm{TV}$ uses $A=1$ risk sets and train units only.}
\label{tab:int_r_occ}
\scriptsize
\begin{tabular}{lcccc}
\toprule
Dataset & Seed & Active frac. & Burstiness & $\mathrm{TV}(\pi^{\mathrm{opp}},\boldsymbol\mu)$\\
\midrule
SensorScope & 30001 & $0.453$ & $0.198$ & $0.309$\\
SensorScope & 30010 & $0.433$ & $0.223$ & $0.300$\\
SensorScope & 30020 & $0.394$ & $0.251$ & $0.271$\\
U-Air & 30001 & $0.495$ & $0.263$ & $0.225$\\
U-Air & 30010 & $0.428$ & $0.190$ & $0.275$\\
U-Air & 30020 & $0.403$ & $0.201$ & $0.205$\\
\bottomrule
\end{tabular}
\end{table}

Table~\ref{tab:int_r_e3} ranks the same methods as
Table~III on those traces ($20$ paired seeds).
Three always-on conclusions survive.  (i)~Four-block
$\mathrm{WorstGroupRMSE}$ remains within the descriptive $+0.5\%$
tolerance versus FedAU-Window
and ObsUse-Window: SensorScope paired relative
change $+0.08\%$ and $+0.12\%$; U-Air $-0.10\%$ and $-0.10\%$.  Holm tests
against those two are not a unique-win claim.  (ii)~Design on for U-Air
still worsens the primaries relative to RAVEN-NoDesign
(paired $+3.65\%$ $\mathrm{RMSE}_{\mu}$ and $+4.07\%$
$\mathrm{WorstGroupRMSE}$; $+9.14\%$ $\Delta_{\mathrm{group}}$), while
slightly improving $\Delta_{c\text{-}s}$ ($-0.26\%$).  The always-on
reversal was about $7\%$ on the primaries; the sign is unchanged.
(iii)~SensorScope $\Delta_{\mathrm{group}}$ remains lower for
RAVEN-Design than for FedAvg (paired $-1.04\%$).

Two always-on headlines do \emph{not} carry over at the same strength.
SensorScope $\Delta_{c\text{-}s}$ versus FedAvg shrinks to
$-0.05\%$ (not distinguishable from zero at conventional $p=0.05$).
Removing Design on SensorScope moves
$\Delta_{\mathrm{group}}$ by only $+0.05\%$ (not significant), so the
always-on ablation statement that Design is the useful SensorScope factor is
not recovered under this occupancy process.  Absolute
$\Delta_{\mathrm{group}}$ also sits higher for every method than in
Table~III (about $1.33$ versus $1.01$ on
SensorScope), which is expected when protocol $R$ is thinned; it is not
a second ranking family and is not mixed into the main table.

\begin{table*}[t]
\centering
\caption{INT-R robustness ranking under intermittent protocol $R$
(mean $\pm$ std over 20 seeds).  Lower is better.  The Design
settings are the same as in Table~III.  Design on for U-Air is also reported.  This is protocol occupancy on station clients,
not mobile-client MCS and not a T-Drive RMSE ranking.}
\label{tab:int_r_e3}
\scriptsize
\setlength{\tabcolsep}{1pt}
\begin{tabular*}{\textwidth}{@{\extracolsep{\fill}}llcccc}
\toprule
Dataset & Method & $\mathrm{RMSE}_{\mu}$ & WorstGroupRMSE &
$\Delta_{\mathrm{group}}$ & $\Delta_{c\text{-}s}$\\
\midrule
\multirow{5}{*}{SensorScope}
& FedAvg-Window     & \shortstack{$7.9373$\\$\pm0.0960$} & \shortstack{$8.6716$\\$\pm0.0976$} & \shortstack{$1.3411$\\$\pm0.0580$} & \shortstack{$1.5866$\\$\pm0.0445$}\\
& TwoStage-H\'ajek    & \shortstack{$7.9362$\\$\pm0.0956$} & \shortstack{$8.6696$\\$\pm0.0957$} & \shortstack{$1.3558$\\$\pm0.0545$} & \shortstack{$1.5941$\\$\pm0.0415$}\\
& FedAU-Window      & \shortstack{$7.9335$\\$\pm0.1036$} & \shortstack{$8.6678$\\$\pm0.1044$} & \shortstack{$1.3540$\\$\pm0.0541$} & \shortstack{$1.5978$\\$\pm0.0425$}\\
& ObsUse-Window     & \shortstack{$7.9305$\\$\pm0.1047$} & \shortstack{$8.6647$\\$\pm0.1054$} & \shortstack{$1.3560$\\$\pm0.0534$} & \shortstack{$1.5983$\\$\pm0.0421$}\\
& RAVEN-Design & \shortstack{$7.9407$\\$\pm0.0957$} & \shortstack{$8.6746$\\$\pm0.0972$} & \shortstack{$\mathbf{1.3272}$\\$\pm0.0607$} & \shortstack{$1.5859$\\$\pm0.0443$}\\
\midrule
\multirow{6}{*}{U-Air}
& FedAvg-Window     & \shortstack{$1.1121$\\$\pm0.0873$} & \shortstack{$1.1682$\\$\pm0.1022$} & \shortstack{$1.3141$\\$\pm0.0679$} & \shortstack{$1.5786$\\$\pm0.0602$}\\
& TwoStage-H\'ajek    & \shortstack{$1.1555$\\$\pm0.1027$} & \shortstack{$1.2179$\\$\pm0.1079$} & \shortstack{$1.4393$\\$\pm0.0613$} & \shortstack{$1.5694$\\$\pm0.0631$}\\
& FedAU-Window      & \shortstack{$1.1126$\\$\pm0.0845$} & \shortstack{$1.1689$\\$\pm0.0985$} & \shortstack{$1.3319$\\$\pm0.0690$} & \shortstack{$1.5951$\\$\pm0.0597$}\\
& ObsUse-Window     & \shortstack{$1.1126$\\$\pm0.0853$} & \shortstack{$1.1689$\\$\pm0.0995$} & \shortstack{$1.3304$\\$\pm0.0693$} & \shortstack{$1.5947$\\$\pm0.0604$}\\
& RAVEN-Design & \shortstack{$1.1528$\\$\pm0.1005$} & \shortstack{$1.2148$\\$\pm0.1040$} & \shortstack{$1.4134$\\$\pm0.0725$} & \shortstack{$\mathbf{1.5632}$\\$\pm0.0663$}\\
& RAVEN-NoDesign & \shortstack{$1.1120$\\$\pm0.0880$} & \shortstack{$1.1680$\\$\pm0.1031$} & \shortstack{$\mathbf{1.2959}$\\$\pm0.0746$} & \shortstack{$1.5672$\\$\pm0.0638$}\\
\bottomrule
\end{tabular*}
\end{table*}

\paragraph{Scope of the supplementary evidence.}
The INT-R block uses the same two station datasets as
Table~III with a different protocol opportunity
process; it does not replace that table.  Cross-vintage Traffic
outputs remain unused.

\section{Support-Hypothesis Audit for A5, A6a, and A6b}
\label{app:a5_a6_audit}

A5, A6a, and A6b are separate hypotheses: the main text separates them so that
the reconstruction results follow from A5 alone, the perturbation form of
Proposition~3 needs A5 together with the bounded interpolation error A6b, and
the sharp form with $\varepsilon_{\mathrm{stab},r}=0$ additionally needs the
stabilizer condition A6a.  This section reports which clauses hold on the
released traces of the formal protocol blocks, reading the frozen
\texttt{EventTraces} only and re-running no training, weighting, or
calibration step.

The traces cover the 20 formal seeds of both datasets, giving 16000
SensorScope and 12800 U-Air client-window records, of which 12346 and 9760
are observed.  Table~\ref{tab:app_a5_a6_audit} reports the decisive reading of
each sub-condition.  The readings substitute the true selection margins
recorded in the trace for the unlogged plug-ins
$\widehat\pi^{\mathrm{opp}}$, $\widehat p^{\mathrm{obs}}$, and
$\widehat q^{\mathrm{use}}$, and use the constants the formal runs actually
resolved: $p_{\min}=q_{\min}=0.05$, $\pi_{\min}=10^{-6}$, $a_{\max}=20$, and
$d_{\max}=10$.  These are the defaults of the window runner
(\texttt{build\_full\_runner}), and the E3 entry point \texttt{run\_e3\_real}
does not override them; $a_{\max}=40$ belongs to the separate E1-R2
weight-safety calibration and to the E2 traffic runner, and is not the value
these blocks ran with.  Every clause below is evaluated at $a_{\max}=20$.
Each reading is
the decisive quantity rather than a mean, because A5, A6a, and A6b bound a
quantity from below or from above and one violating cell settles a clause.

The two selection margins are protocol quantities rather than per-unit
estimates.  Over all $28800$ records the observation rate takes only the two
block-level values $0.1925$ and $0.35$, so it is constant within every
opportunity stratum and contributes no within-stratum variation.  The
usable-update rate is drawn once per client-window and shared by the at most
four risk units of that window, so it reweights client-windows without tilting
within a risk set.  Neither margin varies at the unit level or with the label
on these traces.

\begin{table}[!t]
\centering
\caption{Support-hypothesis audit on the formal-protocol traces.  Each column is
the decisive reading for the dataset; ``active'' restricts to the observed
client-windows entering the estimator.  The ceiling and
$\varepsilon_{\mathrm{stab}}$ rows are read at the true margins on the active
sets of the 20 Design-on runs, with
$\zeta=\pi^{\mathrm{tar}}/\pi^{\mathrm{opp}}$ and against $a_{\max}=20$; the
$\varepsilon_{\mathrm{stab}}$ row is median/p95/max.}
\label{tab:app_a5_a6_audit}
\footnotesize
\setlength{\tabcolsep}{3pt}
\begin{tabular*}{\columnwidth}{@{\extracolsep{\fill}}llcc}
\toprule
 & Sub-condition & SensorScope & U-Air\\
\midrule
A5 & $\min \pi^{\mathrm{opp}}>0$ & $6.33\times10^{-4}$ & $1.39\times10^{-3}$\\
A5 & $\min p^{\mathrm{obs}}\ge p_{\min}$ & $0.1925$ & $0.1925$\\
A5 & $\min q^{\mathrm{use}}>0$ (active) & $0.0055$ & $0.0055$\\
A5 & $\min_r \mathbf{1}^{\top}b>0$ & $5.521$ & $16.89$\\
\midrule
A6a & $q_{\min}$ floor inactive & $5.22\%$ & $5.18\%$\\
A6a & $d_{\max}$ ceiling inactive & $24.22\%$ & $23.75\%$\\
A6a & $\max 1/p^{\mathrm{obs}}$, Design off & $5.195$ & $5.195$\\
A6a & $\max \zeta/p^{\mathrm{obs}}$, Design on & $19.26$ & $\mathbf{23.47}$\\
\midrule
A6a & $\varepsilon_{\mathrm{stab}}$ med/p95/max & $0/0.07/0.55$ & $0/0.05/0.68$\\
\bottomrule
\end{tabular*}
\end{table}

\noindent\textbf{A5 holds.}  The opportunity and observation margins are
strictly positive everywhere in the traces, and the usable-update margin is
strictly positive on every client-window that actually enters the estimator.
The smallest observed usable margin is $5.5\times10^{-3}$, no observed
client-window has a vanishing usable margin, and the smallest window-level
H\'ajek mass is $5.521$ on SensorScope and $16.89$ on U-Air.  The positivity
and normalization steps of Proposition~2 and the denominator in
(24) are therefore not degenerate on these traces.

The one qualification is structural rather than numerical.  A usable-update
margin of exactly zero occurs for $22.84\%$ of SensorScope and $23.75\%$ of
U-Air client-windows.  Those cells are never observed and never usable, so
they contribute neither to a corrected mass nor to any ratio: the count of
zero-margin cells among observed client-windows is $0$ out of $12346$ and $0$
out of $9760$.  This is why A5 is stated on the active support.  Read without
that restriction, the same fact would look like a positivity failure, whereas
its actual content is that the two-stage selection already excludes those
cells from the estimator.

\noindent\textbf{A6a fails; A6b is not tested here.}  The usable-update
clauses are violated.  The
true usable margin falls below the floor $q_{\min}=0.05$ on $5.22\%$ of
observed SensorScope and $5.18\%$ of observed U-Air client-windows, so the
floor is active rather than inactive.  More consequentially,
$1/q^{\mathrm{use}}$ exceeds the ceiling $d_{\max}=10$ on $24.22\%$ and
$23.75\%$ of observed client-windows; over all records, including the
never-observed cells, the share with $q^{\mathrm{use}}<1/d_{\max}$ is
$41.52\%$ and $41.86\%$.

The observation-side clause behaves differently under the two Design settings,
and this is where an earlier reading of this audit was wrong: the clause bounds
the composite weight, not $1/p^{\mathrm{obs}}$.  With Design off the composite
equals $1/p^{\mathrm{obs}}$, whose largest value is $5.195$ because the
observation rate takes only the block-level values $0.1925$ and $0.35$, so the
clause holds on both datasets.  With Design on the composite is
$\zeta/p^{\mathrm{obs}}$ with $\zeta=\pi^{\mathrm{tar}}/\pi^{\mathrm{opp}}$, and
because the opportunity margin ranges down to $6.33\times10^{-4}$ it reaches
$19.26$ on SensorScope and $23.47$ on U-Air, against the ceiling
$a_{\max}=20$ that these runs resolved.  The clause therefore \emph{fails on
U-Air}, on $2.74\%$ of its active client-windows, while it holds on
SensorScope by a narrow margin.  The earlier claim that the reading is
unchanged at $a_{\max}=20$ and at $40$ came from the invalid
$1/p^{\mathrm{obs}}$ form; the Design-on reading is threshold-dependent and
would pass at $40$, which these blocks did not use.

The consequence for the theory is the one recorded in the main text.  The
stabilization error $\varepsilon_{\mathrm{stab}}$ of (37) does not vanish on
these traces: its median is $0$, because most windows leave every stabilizer
inactive at the true margins, but it is strictly positive on $10.7\%$
(SensorScope) and $11.1\%$ (U-Air) of active windows and reaches $0.553$ and
$0.680$ in $\ell_1$ on the reference vector.  Proposition~3 therefore applies
in its $\varepsilon_{\mathrm{stab}}$ form and not in the form that cancels the
error, and the failure is not confined to one layer: the usable-update stage
supplies most of it, and under Design on the opportunity stage supplies the
rest through the composite ceiling.

\noindent\textbf{What this audit does not settle.}  Three limits bound the
reading.  First, A5's clause that the ideal and practical references share
$\mathcal{A}_r$, $M_r$, and $\mathcal{C}_r$ is not recoverable from the
released artifacts: the stored P2 invocation states carry no active-set list,
so that clause is asserted and not measured here.  Second, the reading uses
the true margins in place of the unlogged plug-ins, so it diagnoses the
ideal-side hypothesis rather than the realized clipped weights; the direction
of the plug-in error is not bounded by this audit.  Third, a positive minimum
over 20 fixed seeds is an empirical lower bound and not a population
guarantee, and the same applies to the window masses used for
$\mathbf{1}^{\top}b$.  None of the three limits affects the A5/A6a/A6b split
itself, which is a statement about which hypothesis carries which step of the
proofs, and none of them bears on $\varepsilon_{\mathrm{stab}}$, which is
defined at the true margins by construction.


\section{Coverage-Stress Audit of the Candidate Gate}
\label{app:certificate_audit}

The main paper reports pooled CERTIFIED validity $0.959$ and zero
false-safe openings.  Those rates are empirical.  The adaptive
RESET/CERTIFIED switching is not given a validity proof.
Design is set by configuration, not by this Gate.

\begin{table*}[t]
\centering
\caption{Coverage-stress audit of a candidate aggregate coverage Gate.
Validity is the fraction of CERTIFIED windows with
$C^{\mathrm{LCB}}\le C$.  FSR is the false-safe rate among CERTIFIED
windows with $C<\tau_{\mathrm{cov}}$.  Gate ON is the fraction of all
windows with $\mathrm{Gate}_t=1$.  SensorScope $a=0.40$ and U-Air
$a=0.20$ keep Gate off because the LCB stays below $\tau_{\mathrm{cov}}$.
U-Air $a=1.00$ misses the nominal $95\%$ empirical validity.}
\label{tab:certificate_stress}
\scriptsize
\setlength{\tabcolsep}{3.2pt}
\begin{tabular}{lcccccccc}
\toprule
Dataset & $a$ & $n$ & $n_{\mathrm{cert}}$ & Duty & Validity & FSR & Gate ON & Median $C^{\mathrm{LCB}}$\\
\midrule
\multirow{5}{*}{SensorScope}
& $1.00$ & $62$ & $43$ & $0.694$ & $1.000$ & $0.000$ & $0.694$ & $0.124$\\
& $0.80$ & $62$ & $29$ & $0.468$ & $0.966$ & $0.000$ & $0.452$ & $0.093$\\
& $0.60$ & $62$ & $24$ & $0.387$ & $0.958$ & $0.000$ & $0.306$ & $0.065$\\
& $0.40$ & $62$ & $24$ & $0.387$ & $0.958$ & $0.000$ & $0.000$ & $0.041$\\
& $0.20$ & $62$ & $24$ & $0.387$ & $0.958$ & $0.000$ & $0.000$ & $0.017$\\
\midrule
\multirow{5}{*}{U-Air}
& $1.00$ & $53$ & $8$ & $0.151$ & $0.750$ & $0.000$ & $0.151$ & $0.182$\\
& $0.80$ & $53$ & $11$ & $0.208$ & $0.818$ & $0.000$ & $0.208$ & $0.129$\\
& $0.60$ & $53$ & $15$ & $0.283$ & $0.933$ & $0.000$ & $0.245$ & $0.093$\\
& $0.40$ & $53$ & $31$ & $0.585$ & $0.968$ & $0.000$ & $0.340$ & $0.051$\\
& $0.20$ & $53$ & $34$ & $0.642$ & $1.000$ & $0.000$ & $0.000$ & $0.013$\\
\bottomrule
\end{tabular}
\end{table*}

\begin{figure}[t]
\raggedright
\includegraphics[width=\columnwidth]{fig_certificate_audit.pdf}
\caption{Finite-sample audit of the aggregate coverage certificate.
(a)~CERTIFIED windows: realized coverage $C_t$ versus
$C^{\mathrm{LCB}}_t$.  The dashed line is $y=x$; dotted lines mark
$\tau_{\mathrm{cov}}=0.05$.
(b)~SensorScope $a=0.40$ timeline.  Gray bands are WARMUP; the red mark
is a retained CERTIFIED miscoverage that triggers RESET on the next
window.  The LCB stays below $0.05$, so Gate remains off.}
\label{fig:certificate_audit}
\end{figure}

\section{Target-Coverage Statistic and City Scanner Hybrid Rankings}
\label{app:cityscanner_phi}

Table~\ref{tab:scd_phi_freeze} collects, as a descriptive statistic, the
train-unit zero-target fraction of each block, together with the Design
factor setting used by that block's reported ranking.
The statistic does not select the setting: the Design factor of
(32) is configured, and both of its settings are reported
on the blocks where the trade-off is material.
The statistic does not verify positivity on
target-positive pairs, and none of these counts uses $Y$ or RMSE.
The U-Air EventTrace occurrence share
$13040/51200=25.5\%$ is not a row in this table.
The coverage Gate does not switch Design.

\begin{table}[t]
\centering
\caption{Descriptive target-coverage statistic per block, and the Design
factor setting used by that block's reported ranking.
Zero test: frozen $\pi^{\mathrm{tar}}(k(i),s(i))$ for the occupancy-attached
client; a missing key is $0$.
The statistic is reported for reference and does not select the setting,
and it does not verify positivity on target-positive pairs.
The setting is a configuration, reported together with the ranking it
produced.
U-Air $25.5\%$ is EventTrace occurrence, not this statistic.}
\label{tab:scd_phi_freeze}
\scriptsize
\setlength{\tabcolsep}{3pt}
\begin{tabular*}{\columnwidth}{@{\extracolsep{\fill}}lcccc}
\toprule
Block & $n_{\mathrm{train}}$ & $n_{\mathrm{zero}}$ & coverage & Design used\\
\midrule
SensorScope (Table~III) & $10285$ & $0$ & $0$ & on\\
U-Air (Table~III) & $5688$ & $2016$ & $35.4\%$ & off\\
T-Drive hybrid (Table~II) & $285659$ & $135268$ & $47.4\%$ & off\\
NSW-Traffic four-region (Table~\ref{tab:traffic_four_region}) & $18576$ & $0$ & $0$ & on\\
City Scanner NYC {PM}$_{2.5}$ (Table~\ref{tab:cityscanner_hybrid}) & $9343$ & $0$ & $0$ & on\\
City Scanner P2 meas.\ $O$ (Table~\ref{tab:cityscanner_p2_measured_o}) & $10141$ & $0$ & $0$ & on\\
\bottomrule
\end{tabular*}
\end{table}

SensorScope and U-Air look up test-split client--stratum mass
$\pi^{\mathrm{tar}}_{k,s}$.  T-Drive looks up test-split stratum
frequency.  City Scanner blocks assign the test-group share $\mu_{h(i)}$
to training atoms, with all four regions positive.  NSW-Traffic uses a
four-region freeze.  These are distinct lookup objects.  The same
coverage statistic on one map does not mean that the SensorScope/U-Air
support statistic transferred.  The City Scanner value follows from
positive regional $\mu$, not from checking target-positive pairs.

\subsection{City Scanner New York {PM}$_{2.5}$}
\label{app:cityscanner_nyc}

This block uses Design on with a locally constructed
public drive-by
{PM}$_{2.5}$ map~\cite{wang2023cityscanner}.
The source files are the New York Pilot~1 and Pilot~2 {PM} CSVs on
Zenodo~\footnote{\url{https://doi.org/10.5281/zenodo.7961851}}
({CC BY} 4.0).  Boston, Beirut, and {NO}$_2$ files are not used.
$Y$ is the mean of finite calibrated {PM}$_{2.5}$ ($\mu$g/m$^3$) inside
an occupied $500$\,m $\times$ $30$\,min cell.
GPS places visits only.
Federated clients are occupied grid cells, not vehicles.
Physical filters drop non-finite latitude, longitude, time, or
\texttt{pm25}, and drop $\texttt{pm25}<0$ or $\texttt{pm25}>1000$.
Empty cell-slots are absent, not imputed.

The public target $h(i)$ is four quadrants from the median cell indices
$(g_x,g_y)$ over all occupied cells, without $Y$.
$\boldsymbol\mu$ is the test-split region share, without $Y$.
Opportunity strata are $\mathrm{region}{::}\mathrm{block}{::}\mathrm{weekday}$
(at most $32$).
$R$ is real visit occupancy: a client occupies a unit if and only if
that cell has at least one kept sample in that slot.
$O$ and $U$ remain Complete-aligned controls
($p_{\mathrm{obs}}=0.35$, $S_{\max}=5$).
The ranking family is therefore hybrid, not end-to-end MCS.

Processed counts: $14723$ atomic units, $446$ occupied cells, $2776$
occupied slots, and train/validation/test $9343/3290/2090$.
All four regions have test units
($496$, $992$, $148$, $454$).
Table~\ref{tab:scd_phi_freeze} records the target-coverage statistic
($0$ of $9343$); Design is on for this block.
Shared EventTraces use seeds $32001$--$32010$ and $100$ windows obtained
by chronological equal-count bins of the $1666$ train $30$\,min slots.
The coverage Gate is unused.

Table~\ref{tab:cityscanner_hybrid} ranks the same windowed family.
$\mathrm{WorstGroupRMSE}$ is the maximum error over the four quadrants.
RAVEN-Design is lower than FedAU-Window and ObsUse-Window on
paired $\mathrm{RMSE}_{\mu}$ (mean relative $-1.71\%$ and $-1.34\%$)
and on worst-group RMSE ($-1.56\%$ and $-1.23\%$).
FedAvg-Window remains lowest on both prediction endpoints.
RAVEN-Design is smallest on $\Delta_{\mathrm{group}}$.
The block is not mixed into Table~III or
Table~II, and it is not a confirmation of those
tables.
A separate Pilot~2 ranking with measured observation is
Table~\ref{tab:cityscanner_p2_measured_o}.

\begin{table}[t]
\centering
\caption{City Scanner New York {PM}$_{2.5}$ hybrid ranking
(mean $\pm$ std, 10 seeds, $100$ windows).  Lower is better.
FedAvg-Window is lowest on $\mathrm{RMSE}_{\mu}$ and WorstGroupRMSE;
RAVEN-Design is lower than FedAU-Window and ObsUse-Window on
those two and smallest on $\Delta_{\mathrm{group}}$.}
\label{tab:cityscanner_hybrid}
\scriptsize
\setlength{\tabcolsep}{2.5pt}
\begin{tabular}{@{}lcccc@{}}
\toprule
Method & $\mathrm{RMSE}_{\mu}$ & WorstGroupRMSE &
$\Delta_{\mathrm{group}}$ & $\Delta_{c\text{-}s}$\\
\midrule
FedAvg-Window & \shortstack{$\mathbf{0.2379}$\\$\pm0.0839$} & \shortstack{$\mathbf{0.2489}$\\$\pm0.0856$} & \shortstack{$0.1953$\\$\pm0.0342$} & \shortstack{$0.8123$\\$\pm0.0191$}\\
TwoStage-H\'ajek & \shortstack{$0.2510$\\$\pm0.1044$} & \shortstack{$0.2629$\\$\pm0.1068$} & \shortstack{$0.0952$\\$\pm0.0307$} & \shortstack{$\mathbf{0.5305}$\\$\pm0.0286$}\\
FedAU-Window & \shortstack{$0.2582$\\$\pm0.1028$} & \shortstack{$0.2700$\\$\pm0.1048$} & \shortstack{$0.3216$\\$\pm0.0412$} & \shortstack{$1.1403$\\$\pm0.0538$}\\
ObsUse-Window & \shortstack{$0.2566$\\$\pm0.1004$} & \shortstack{$0.2685$\\$\pm0.1023$} & \shortstack{$0.3187$\\$\pm0.0414$} & \shortstack{$1.1369$\\$\pm0.0547$}\\
RAVEN-Design & \shortstack{$0.2549$\\$\pm0.1057$} & \shortstack{$0.2671$\\$\pm0.1079$} & \shortstack{$\mathbf{0.0706}$\\$\pm0.0266$} & \shortstack{$0.5676$\\$\pm0.0311$}\\
\bottomrule
\end{tabular}

\begin{minipage}{\columnwidth}
\raggedright\scriptsize
$R$ is real visit occupancy; $O$/$U$ are Complete-aligned
($p_{\mathrm{obs}}=0.35$, $S_{\max}=5$).
WorstGroupRMSE is the maximum over four coordinate quadrants.
RAVEN-Design is used (Design on).
Clients are occupied cells, not vehicles.
Not mixed into Table~III or
Table~II.
\end{minipage}
\end{table}

\subsection{City Scanner Pilot~2 with measured observation}
\label{app:cityscanner_p2_measured_o}

This block keeps the same Design setting and the same windowed family, but
changes how observation is realized.
The source files are the New York Pilot~2 {PM} and {NO}$_2$ CSVs on
Zenodo~\footnote{\url{https://doi.org/10.5281/zenodo.7961851}}
({CC BY} 4.0).  Pilot~1, Boston, and Beirut are not used.
$Y$ is the mean of finite calibrated {PM}$_{2.5}$ ($\mu$g/m$^3$) inside
an occupied $500$\,m $\times$ $30$\,min cell.
Federated clients are occupied grid cells, not vehicles.

$R$ is real visit occupancy on the union of kept-{PM} GPS and {NO}$_2$ GPS
at the same grid: a client occupies a unit if and only if that cell has at
least one kept GPS sample in that slot.
$O=1$ if and only if the cell-slot contains finite {PM}$_{2.5}$ in
$[0,1000]$; this is not a Bernoulli $p_{\mathrm{obs}}$ control.
{NO}$_2$ contributes GPS occupancy only and is not the completion target.
$U$ remains a Complete-aligned control ($S_{\max}=5$).
The ranking family is therefore hybrid, not an end-to-end MCS system and
not a jointly measured observation-and-network process.

Occupancy is $14960$ cell-slots, $416$ cells, and $2402$ slots
($10141/2835/1984$ train/val/test), with measured $O=0$ on $1931$ slots
($12.9\%$) and $13029$ labeled $O=1$ units.
Region test occupancy is $1289$, $268$, $144$, and $283$.
RAVEN-Design is used (Design on); seeds $33001$--$33010$ and
$100$ windows from $1441$ train slots.
Paired relative RMSE versus FedAU-Window, ObsUse-Window, and FedAvg-Window is
$-0.43\%$, $-0.45\%$, and $-0.65\%$ (worst-group $-0.44\%$, $-0.47\%$,
$-0.65\%$).
Not mixed into Table~III,
Table~II, or Table~\ref{tab:cityscanner_hybrid},
and not a confirmation of those tables.

\begin{table*}[t]
\centering
\caption{City Scanner New York Pilot~2 ranking with measured observation
(mean $\pm$ std, 10 seeds, $100$ windows).  Lower is better.
TwoStage-H\'ajek is lowest on $\mathrm{RMSE}_{\mu}$ and WorstGroupRMSE;
RAVEN-Design is smallest on $\Delta_{\mathrm{group}}$ and
lower than FedAU-Window, ObsUse-Window, and FedAvg-Window on the two prediction
endpoints.}
\label{tab:cityscanner_p2_measured_o}
\scriptsize
\setlength{\tabcolsep}{2.5pt}
\begin{tabular*}{\textwidth}{@{\extracolsep{\fill}}lcccc}
\toprule
Method & $\mathrm{RMSE}_{\mu}$ & WorstGroupRMSE &
$\Delta_{\mathrm{group}}$ & $\Delta_{c\text{-}s}$\\
\midrule
FedAvg-Window & $0.7182{\pm}0.0208$ & $0.7359{\pm}0.0204$ & $0.0769{\pm}0.0231$ & $0.7066{\pm}0.0162$\\
TwoStage-H\'ajek & $\mathbf{0.7097}{\pm}0.0204$ & $\mathbf{0.7273}{\pm}0.0200$ & $0.0607{\pm}0.0206$ & $\mathbf{0.3893}{\pm}0.0265$\\
FedAU-Window & $0.7173{\pm}0.0289$ & $0.7351{\pm}0.0291$ & $0.2283{\pm}0.0433$ & $1.0034{\pm}0.0383$\\
ObsUse-Window & $0.7175{\pm}0.0288$ & $0.7352{\pm}0.0290$ & $0.2314{\pm}0.0429$ & $1.0059{\pm}0.0380$\\
RAVEN-Design & $0.7135{\pm}0.0206$ & $0.7311{\pm}0.0200$ & $\mathbf{0.0465}{\pm}0.0119$ & $0.4172{\pm}0.0214$\\
\bottomrule
\end{tabular*}

\par\vspace{2pt}
\noindent\raggedright\scriptsize
$R$ is real visit occupancy on union GPS; $O$ is measured from finite
{PM}$_{2.5}$; $U$ is Complete-aligned ($S_{\max}=5$).
WorstGroupRMSE is the maximum over four coordinate quadrants.
RAVEN-Design is used (Design on).
Clients are occupied cells, not vehicles.
Not mixed into Table~III,
Table~II, or Table~\ref{tab:cityscanner_hybrid}.
\end{table*}


